% DOCUMENT322 draft: grouped growing marginal order and critical moment rigidity; new exact-PDF assessment pending.
\documentclass[11pt]{article}
\usepackage[margin=1in]{geometry}
\usepackage[T1]{fontenc}
\usepackage{lmodern,amsmath,amssymb,amsthm,microtype}
\usepackage[colorlinks=true,allcolors=blue]{hyperref}
\pdfinfoomitdate=1
\pdftrailerid{}
\ifdefined\pdfsuppressptexinfo\pdfsuppressptexinfo=15\fi
\newtheorem{theorem}{Theorem}
\newtheorem{lemma}[theorem]{Lemma}
\newtheorem{proposition}[theorem]{Proposition}
\newtheorem{corollary}[theorem]{Corollary}
\newtheorem{remark}[theorem]{Remark}
\newcommand{\F}{\mathbb F}
\newcommand{\R}{\mathbb R}
\newcommand{\diag}{\operatorname{diag}}
\newcommand{\gap}{\operatorname{gap}}
\newcommand{\wt}{\operatorname{wt}}
\newcommand{\Unif}{\operatorname{Unif}}
\providecommand{\SourceCommit}{SOURCE-NOT-BUILT}
\title{Product Marginals and Density-Moment Thresholds\\for Ising Transfer Kernels}
\author{Lluis Eriksson\\\small Independent researcher}
\date{13 November 2026}
\begin{document}
\maketitle
\begin{abstract}
Fix a positive Ising coupling. We construct strictly positive Boolean
weights whose labelled marginals on every set of at most a fixed positive
fraction of coordinates equal uniform product marginals, while exponentially
many normalized transfer eigenvalues approach one exponentially fast.
The Perron stationary law has the same marginals. Binary codes with positive
primal and dual distances, a positive product admixture, and finite block
estimates supply the construction. An exact parity spectrum distinguishes
transfer relaxation from single-site heat-bath relaxation.
Conversely, for Boolean noise parameter $0<\rho<1$, density moments of order
$s>(1+\rho)/(1-\rho)$ imply a dimension-independent normalized transfer gap,
even allowing zero weights. Below this exponent, every norm budget $B>1$
admits positive densities with exponentially collapsing gaps.
At the critical exponent $s=(1+\rho)/(1-\rho)$ the exact worst-case norm
boundary is $B_*=2^{1/s}$: smaller budgets give a positive gap, whereas
at $B_*$ and above collapse is possible with any prescribed density floor
below one. We also identify the sharp worst-case boundary for the
asymptotic escaped critical moment
$\Theta=\lim_{K\to\infty}\sup_L\mathbb E[r_L^s\mathbf1_{r_L>K}]$.
The condition $\Theta<1$ gives an explicit uniform gap, while a collapsing
family has $\Theta=1$. More strongly, a single collapsing family has tail moment above one
at most one and total moment at most two in every dimension; it meets
the same tail cap simultaneously at every larger cutoff. For fixed
$0\le K<1$, this tail cap instead guarantees a positive gap. Nevertheless stable one-vertex spikes realize every
finite $\Theta>0$, so critical-power uniform integrability is sufficient
but not necessary. These conclusions use classical hypercontractivity,
sign compression, balanced product tilts and rank-one perturbations.
For $k_s=(s-1)/s<K<1$, finite parity backgrounds also give collapse
at the attained tail budget $1+C_s(1-K)K^{s-1}<2$,
with density floor exactly at least $K$, where
$C_s=s^s/(s-1)^{s-1}$. We prove that this is the exact universal
strict-tail threshold: every smaller budget yields an explicit uniform
gap for arbitrary nonnegative densities, without a density floor.
The proof combines reverse hypercontractivity, normalized sign compression
and truncation; a sharp scalar total-moment envelope supplies a simpler
intermediate bound. We additionally connect exact weighting-law marginals to the critical moment
budget. For every $0<\alpha<1/2$, positive densities with prescribed floor
below one retain exactly uniform marginals through $\lfloor L^\alpha\rfloor$
and norm at most $2^{1/s}$ while their first two eigenvalues tend to one.
Conversely, at any bounded critical norm, matching through every fixed
positive linear fraction forces operator-norm convergence to the unweighted
noise operator and normalized gap $1-\rho$. This implication uses the published
homogeneous moment bound of Kirshner and Samorodnitsky. Critical-power uniform
integrability supplies a separate rigidity criterion at any diverging marginal
order. The intervening sublinear threshold remains unresolved.
External priority is unestablished. Global many-body weights are allowed;
no fixed-range Gibbs, local-dynamics or physical-gap conclusion follows.
\end{abstract}


Section~\ref{sec:growing-marginals} connects exact weighting-law marginal matching and the critical moment budget in a single family, and proves rigidity for positive linear matching fractions.

The constructions serve different conclusions. The code family in
Theorem~\ref{thm:main} preserves labelled marginals and produces many slow
transfer modes. The product tilts in Theorems~\ref{thm:tilts}
and~\ref{thm:critical-exact} establish the moment-budget boundary, without
the code family's marginal matching. The critical-tail section gives a sufficient
tail condition, stable counterexamples to its necessity, and collapse at
the finite-cutoff endpoint in Theorem~\ref{thm:finite-tail-endpoint}.
Section~\ref{sec:background} then optimizes finite parity backgrounds
to extend the sub-two tail-budget collapse region below cutoff one.
Section~\ref{sec:attained-cutoff} attains the constructive budget
and allows the floor to equal the cutoff. Theorem~\ref{thm:scalar-envelope}
gives the exact scalar total-moment envelope. Theorem~\ref{thm:universal-tail}
then proves necessity of the constructed tail budget for every collapsing
family. Readers interested in this sharp tail classification can start at
Sections~\ref{sec:scalar-envelope}--\ref{sec:universal-tail}, using
Theorem~\ref{thm:attained-cutoff} for the matching endpoint.
These are separate families; their properties are not asserted jointly.

\section{The question and the transfer convention}
Local marginal agreement and rapid relaxation are different kinds of
information. The distinction is familiar in limited independence and
Markov-chain theory, but a comparison must specify the dynamics: a slow
single-site heat-bath chain does not automatically give a slow transfer
operator. We make this distinction within a single fixed-coupling transfer
family and require exact, rather than approximate, agreement of marginal
input.

Let $\Omega_L=\F_2^L$, with Hamming distance $d(x,y)$. Fix $\beta>0$ and put
$q=e^{-2\beta}$. Our unweighted kernel is
\begin{equation}\label{eq:kernel}
 M_\beta(x,y)=q^{d(x,y)}
 =e^{-\beta L}\exp\!\left(\beta\sum_{i=1}^L(-1)^{x_i+y_i}\right).
\end{equation}
For a strictly positive probability law $\mu$ on $\Omega_L$, define
\begin{equation}\label{eq:weighted}
 A_\mu=\diag(\sqrt\mu)M_\beta\diag(\sqrt\mu).
\end{equation}
The kernel $M_\beta$ is a tensor power of
$\left(\begin{smallmatrix}1&q\\q&1\end{smallmatrix}\right)$ and is positive
definite. The same is true of $A_\mu$, which also has strictly positive
entries. Write its eigenvalues in decreasing order as
$\lambda_1>\lambda_2\ge\cdots>0$, counting multiplicities, and set
\[
 \gap(A_\mu)=1-\lambda_2/\lambda_1.
\]
Perron--Frobenius strictness gives a positive gap at each finite $L$; a
bound uniform in $L$ is a separate matter. Multiplying either kernel by
a positive scalar does not change any eigenvalue ratio.

Let $u_L=\Unif(\Omega_L)$. Tensorization gives
\begin{equation}\label{eq:product}
 \lambda_2(A_{u_L})/\lambda_1(A_{u_L})
       =\frac{1-q}{1+q}=\tanh\beta.
\end{equation}
For a law $\nu$, the literal marginal input of order $r$ is
\begin{equation}\label{eq:input}
 \mathcal M_r(\nu)=(\nu_S:S\subseteq[L],\ |S|\le r),
\end{equation}
where coordinate labels are retained. Full conditional distributions,
entrywise weight bounds, and the transfer matrix itself are not included
in this tuple.

\begin{theorem}\label{thm:main}
For each fixed $\beta>0$ there are constants $\alpha,R,\gamma>0$ such
that, for all sufficiently large integers $L$, a strictly positive law
$\mu_L$ on $\Omega_L$ has the following properties. Put
$r_L=\lfloor\alpha L\rfloor$ and $m_L=2^{\lceil RL\rceil}$.
\begin{enumerate}
\item $\mathcal M_{r_L}(\mu_L)=\mathcal M_{r_L}(u_L)$.
\item For every $2\le j\le m_L$,
\begin{equation}\label{eq:slow}
 0<1-\lambda_j(A_{\mu_L})/\lambda_1(A_{\mu_L})
                 \le 7e^{-\gamma L}.
\end{equation}
\item Let $v>0$ be the unit Euclidean Perron vector of $A_{\mu_L}$.
The reversible Perron Doob chain
\begin{equation}\label{eq:doob}
 P_L(x,y)=\frac{A_{\mu_L}(x,y)v_y}{\lambda_1 v_x},
 \qquad \pi_L(x)=v_x^2,
\end{equation}
has $\mathcal M_{r_L}(\pi_L)=\mathcal M_{r_L}(u_L)$ and the same
normalized spectrum as $A_{\mu_L}$.
\end{enumerate}
One admissible choice is
\begin{equation}\label{eq:parameters}
 R=\min\left\{\frac{1-h(1/4)}2,\frac{\beta}{4\log 2}\right\},
 \quad h(\alpha)<R/2,\quad
 \gamma=\beta/2-R\log 2\ge\beta/4,
\end{equation}
with $0<\alpha<1/2$ and binary entropy
$h(p)=-p\log_2p-(1-p)\log_2(1-p)$.
\end{theorem}

There are at least $m_L-1$ slow nonleading modes in this statement.
We do not assert that their eigenvalues coincide. The positivity and
strict gap at each finite size are compatible with their exponentially
small separation from the Perron root. The theorem does not impose a
uniform lower bound on the weights or a bounded interaction range.

\section{Codes with two distance conditions}
We include the coding argument to make all hypotheses of the spectral
construction explicit. Simultaneously good primal and dual codes are
classical; stronger algorithmic constructions are available
\cite{BrehmResch}. The relation between code projections and limited
independence is standard background for the marginal comparison
\cite{AGM}.

\begin{lemma}\label{lem:code}
Suppose $0<\alpha,\delta<1/2$ and
$h(\alpha)<R<1-h(\delta)$. Put $k=\lceil RL\rceil$.
If $k<L$ and
\begin{equation}\label{eq:codegate}
 2^{k-L(1-h(\delta))}+2^{-k+Lh(\alpha)}<1,
\end{equation}
there is a binary linear code $C\subseteq\F_2^L$ of dimension $k$
such that every nonzero codeword has weight at least $\delta L$ and
$C^\perp$ has no nonzero word of weight at most $\lfloor\alpha L\rfloor$.
Every coordinate projection of $C$ onto at most $\lfloor\alpha L\rfloor$
coordinates is surjective. The displayed conditions hold for all
sufficiently large $L$.
\end{lemma}
\begin{proof}
Choose a uniformly random binary $k\times L$ matrix $G$. For each fixed
nonzero row vector $z\in\F_2^k$, the word $zG$ is uniform in $\F_2^L$.
The probability that some such $zG$ has weight less than $\delta L$ is
at most
\[
 2^k2^{-L}\sum_{j\le\delta L}\binom Lj
       \le 2^{k-L(1-h(\delta))}.
\]
Failure of full row rank is included in this event, since it produces
a zero word from a nonzero message. For a fixed nonzero column vector
$w\in\F_2^L$, the vector $Gw$ is uniform in $\F_2^k$. Thus the probability
of a nonzero dual word with weight at most $\lfloor\alpha L\rfloor$ is
at most
\[
 2^{-k}\sum_{1\le j\le\alpha L}\binom Lj
       \le 2^{-k+Lh(\alpha)}.
\]
Their sum is less than one, so both failures can be avoided simultaneously.

For the entropy estimate, if $0<p<1/2$, set $z=p/(1-p)<1$. Since
$z^j\ge z^{pL}$ for $j\le pL$, the binomial theorem gives
\[
 \sum_{j\le pL}\binom Lj\le z^{-pL}(1+z)^L=2^{Lh(p)}.
\]
This bound applies to the inclusive sums and hence to the strict-weight
failure event as well. No integrality assumption on $\delta L$ is needed.

If the columns indexed by $S$ were dependent, their relation would give
a nonzero word of $C^\perp$ supported in $S$. Consequently $G_S$ has
rank $|S|$ whenever $|S|\le\lfloor\alpha L\rfloor$. Uniform messages
therefore give uniform projections onto $\F_2^S$. Finally, the exponents
in~\eqref{eq:codegate} tend to minus infinity, since $k=RL+O(1)$ and
both entropy inequalities are strict.
\end{proof}

\section{A finite positive concentration estimate}
The spectral step needs only the size and minimum distance of the code.
The dual distance will enter through the marginal statement.

\begin{proposition}\label{prop:finite}
Let $C\subseteq\F_2^L$ be a binary linear code with size $m\ge2$ and
minimum nonzero weight $d_C$. Suppose its coordinate projections onto
sets of size at most $r$ are surjective. For $0<\eta\le1/2$, let
\begin{equation}\label{eq:mixture}
 \mu=(1-\eta)\Unif(C)+\eta u_L,
 \qquad s=(m-1)e^{-2\beta d_C}.
\end{equation}
Then $\mu$ is strictly positive, $\mathcal M_r(\mu)=\mathcal M_r(u_L)$,
and every $2\le j\le m$ satisfies
\begin{equation}\label{eq:finite}
 0<1-\lambda_j(A_\mu)/\lambda_1(A_\mu)
       \le 2s+2m(\sqrt\eta+\eta).
\end{equation}
\end{proposition}
\begin{proof}
Uniform codewords have uniform projections by surjectivity, and a mixture
of two laws with the same marginals has those same marginals. The positive
product part gives strict positivity at every configuration.

Every codeword has probability
\[
 a=(1-\eta)/m+\eta/2^L\ge1/(2m).
\]
The principal block of $A_\mu$ indexed by $C$ is $aM_{CC}$. Its diagonal
is $a$, and the off-diagonal row sums of $M_{CC}$ are at most $s$:
distinct codewords differ by a nonzero codeword of weight at least $d_C$.
Thus all eigenvalues of $M_{CC}$ lie in $[1-s,1+s]$. For example, the
norm of its symmetric off-diagonal part is bounded by the maximum
absolute row sum. Principal interlacing gives
\begin{equation}\label{eq:interlace}
 \lambda_j(A_\mu)\ge\lambda_m(A_\mu)\ge a(1-s)
       \quad(1\le j\le m).
\end{equation}
Independently, a coordinate Rayleigh vector gives $\lambda_1(A_\mu)\ge a$.

Write $A_\mu$ in $C,C^c$ blocks. The outside principal block is positive
semidefinite, and its trace is
$\sum_{C^c}\mu=\eta(1-m/2^L)\le\eta$, so its norm is at most $\eta$.
For the rectangular cross block $B$, all entries of $M_\beta$ are at most
one, and therefore
\[
 \|B\|^2\le\|B\|_{\rm F}^2
 \le\left(\sum_C\mu\right)\left(\sum_{C^c}\mu\right)\le\eta.
\]
The symmetric block matrix with off-diagonal blocks $B,B^*$ and zero
diagonal has norm $\|B\|$. Consequently
\[
 \bigl\|A_\mu-\diag(aM_{CC},0)\bigr\|\le\sqrt\eta+\eta,
 \qquad
 \lambda_1(A_\mu)\le a(1+s)+\sqrt\eta+\eta.
\]
Subtract~\eqref{eq:interlace}, divide by $\lambda_1\ge a$, and use
$a\ge1/(2m)$ to obtain~\eqref{eq:finite}. Perron strictness proves the
strict lower inequality. The upper bound is valid even when it is
larger than one; no small-size success is claimed.
\end{proof}

\begin{proof}[Proof of Theorem~\ref{thm:main}, parts one and two]
Use~\eqref{eq:parameters} and Lemma~\ref{lem:code} with $\delta=1/4$.
The choice of $R$ satisfies $0<R<1-h(1/4)$ and $R\log2\le\beta/4$.
Choose codes for sufficiently large $L$, and put
\[
 m=2^{\lceil RL\rceil},\qquad \eta=e^{-2\beta L}/m^2.
\]
This number lies in $(0,1/4]$. Since $d_C\ge L/4$ and
$m\le2e^{R(\log2)L}$,
\[
 2(m-1)e^{-2\beta d_C}\le4e^{-(\beta/2-R\log2)L}
          =4e^{-\gamma L}.
\]
Also, using $m\ge2$ and $\gamma\le\beta$,
\[
 2m(\sqrt\eta+\eta)
   =2e^{-\beta L}+2e^{-2\beta L}/m
   \le3e^{-\beta L}\le3e^{-\gamma L}.
\]
Proposition~\ref{prop:finite} now proves both claims. The condition $k<L$
ensures $m<2^L$, so every asserted eigenvalue index exists.
\end{proof}

\section{Stationary marginals survive the Perron transform}
The weighting law in a transfer congruence need not equal the stationary
law of its associated Markov chain. Matching the former alone would leave
a significant interpretive gap. Code symmetry pays the latter exactly.

\begin{lemma}\label{lem:symmetry}
For the mixture~\eqref{eq:mixture}, the Perron stationary law
$\pi(x)=v_x^2$ satisfies $\mathcal M_r(\pi)=\mathcal M_r(u_L)$.
\end{lemma}
\begin{proof}
For every $c\in C$, translation $x\mapsto x+c$ preserves $\mu$, Hamming
distance, and hence $A_\mu$. The corresponding permutation commutes with
$A_\mu$ and takes a positive unit Perron vector to another such vector.
Uniqueness forces $v_{x+c}=v_x$, so $\pi$ is invariant under translations
by $C$. If $|S|\le r$, every translation on $\F_2^S$ lifts to some
$c\in C$ by the assumed surjectivity. The $S$-marginal of $\pi$ is
therefore invariant under all translations and must be uniform.
\end{proof}

\begin{proof}[Completion of Theorem~\ref{thm:main}]
The Perron eigenvector equation gives row sums one for~\eqref{eq:doob}.
Symmetry gives detailed balance:
\[
 \pi(x)P_L(x,y)=v_xA_\mu(x,y)v_y/\lambda_1
              =\pi(y)P_L(y,x).
\]
Conjugation by $\diag(v)$ shows that $P_L$ is similar to $A_\mu/\lambda_1$.
It has strictly positive entries and the claimed spectrum. Apply
Lemma~\ref{lem:symmetry} to obtain its stationary marginal assertion.
For the uniform weight, the Perron vector and stationary law are uniform,
and~\eqref{eq:product} supplies the comparison chain.
\end{proof}

\section{An exact parity control with different dynamics}
The positive parity family
\begin{equation}\label{eq:parity}
 \mu_t(x)=2^{-L}(1+t\chi_{[L]}(x)),\qquad |t|<1,
 \qquad \chi_S(x)=(-1)^{\sum_{i\in S}x_i},
\end{equation}
has all proper marginals equal to product marginals. Its single-site
heat-bath gap tends to zero as $|t|\to1$ at fixed size $L\ge2$ and fixed
finite positive site rates~\cite{ErikssonParity}. Nevertheless, its transfer ratio at fixed
$\beta$ is bounded uniformly away from one, as follows.

\begin{proposition}\label{prop:parity}
Let $L\ge2$, $\rho=\tanh\beta$, and $|t|<1$. Up to a common positive
scalar, the spectrum of $A_{\mu_t}$ is the union over unordered complement
pairs $\{S,S^c\}$ of the two roots
\begin{equation}\label{eq:parityroots}
 \frac{z_s\pm\sqrt{z_s^2-4\Delta}}2,
 \quad z_s=\rho^s+\rho^{L-s},\quad s=|S|,
 \quad \Delta=(1-t^2)\rho^L.
\end{equation}
The Perron root is the upper root at $s=0$, and the second eigenvalue is
the upper root at $s=1$. In particular,
\begin{equation}\label{eq:paritybound}
 \frac{\lambda_2(A_{\mu_t})}{\lambda_1(A_{\mu_t})}
 \le\frac{\rho+\rho^{L-1}}{1+\rho^L}
 \le\tanh(2\beta)<1.
\end{equation}
The last bound is sharp as a supremum over $L\ge2$ and $|t|<1$.
\end{proposition}
\begin{proof}
The row-normalized product transfer $P_\beta=M_\beta/(1+q)^L$ acts on
$\chi_S$ by $\rho^{|S|}$. Write $D_t=\diag\sqrt{1+t\chi_{[L]}}$.
The weighted congruence, apart from a common scalar, is $D_tP_\beta D_t$,
which is similar to $P_\beta D_t^2$ since $D_t$ is invertible.
Multiplication by $\chi_{[L]}$ exchanges $\chi_S$ and $\chi_{S^c}$.
The coefficient block on that ordered pair is
\[
 \begin{pmatrix}
 \rho^s&t\rho^s\\ t\rho^{L-s}&\rho^{L-s}
 \end{pmatrix}.
\]
Its trace is $z_s$ and its determinant is $\Delta$, yielding
\eqref{eq:parityroots}. The $2^{L-1}$ complement pairs exhaust the
character basis, with repetitions counted as usual.

The largest $z_s$ is $z_0=1+\rho^L$, and the largest among proper pairs
is $z_1=\rho+\rho^{L-1}$. The upper root increases with $z_s$.
The lower root in the empty/full pair is at most $\rho^L$, whereas
the upper root for $s=1$ is at least $\rho$. Every other lower root
is no larger than its own upper root. This identifies the two largest
eigenvalues.

For fixed $\Delta>0$, the upper root divided by $z$ is
\[
 \frac{1+\sqrt{1-4\Delta/z^2}}2,
\]
which increases with $z\ge2\sqrt\Delta$. Comparing $z_1\le z_0$ proves
the first inequality in~\eqref{eq:paritybound}. For the second, write
$b=\rho^{L-2}\le1$; the ratio is $\rho(1+b)/(1+\rho^2b)$, increasing
in $b$. Its maximum is $2\rho/(1+\rho^2)=\tanh(2\beta)$.
For $L=2$ and $t\to1$, $\Delta\to0$, so equality is approached.
\end{proof}

This control rules out a tempting shortcut: the known parity heat-bath
obstruction cannot be transferred by changing the name of the operator.
The code construction uses a different geometric ingredient, a set of
exponentially many mutually separated configurations while retaining
surjective coordinate projections.

\section{What marginal certificates cannot establish}
\begin{corollary}\label{cor:certificate}
Fix the constants of Theorem~\ref{thm:main}. Suppose $F$ is a valid lower
bound on the transfer gap for every strictly positive weight, and factors
through the tuple
\[
 (L,\beta,\mathcal M_{r_L}(\mu),\mathcal M_{r_L}(\pi)).
\]
At the product input and for all sufficiently large $L$, its value is at
most $7e^{-\gamma L}$. The same conclusion holds if only one of the two
marginal tuples is supplied.
\end{corollary}
\begin{proof}
At the same $L$ and $\beta$, the constructed weight and its stationary
law supply exactly the same input as the uniform weight and uniform
stationary law. A function of that input takes the same value on both.
Validity on the constructed example and~\eqref{eq:slow} for $j=2$
force the asserted bound.
\end{proof}

The corollary does not assert that every such functional is zero or that
the product chain is slow. Its actual gap is $1-\tanh\beta$. Nor does
it concern a certificate using full conditionals, uniform density bounds,
the operator entries, or code membership: those are additional information.
In particular, the weights here can have exponentially small probabilities
and global many-body interactions. The theorem says nothing against a
result restricted to uniformly controlled, fixed-range Gibbs weights.
The Perron Doob chain updates a full configuration and is not a local
single-site dynamics. We make no claim about continuum Yang--Mills,
a physical mass gap, entropy tensorization, or a specified spatial mixing
rate. The many-mode conclusion is spectral; no separate theorem about
metastable sets is implicit.

\section{Upper density control without a lower bound}\label{sec:density}
The obstruction above allows the maximum density relative to product to
increase with dimension. A uniform upper bound changes the conclusion;
a positive lower bound is not required. Write $\rho=\tanh\beta$ and let
$P_\rho$ be the Boolean noise operator in $L^2(u_L)$, whose Walsh
multipliers are $\rho^{|S|}$. For a nonnegative density $r$ with
$\mathbb E_{u_L}r=1$, put
\[
 T_r=\diag(\sqrt r)P_\rho\diag(\sqrt r).
\]
Since $P_\rho=M_\beta/(1+e^{-2\beta})^L$, for $\mu=2^{-L}r$ we have
$A_\mu=2^{-L}(1+e^{-2\beta})^L T_r$. Thus the eigenvalue ratios agree.
Here zeros of $r$ are allowed; we keep the zero eigenvalues outside its
support in the full $2^L$-dimensional list.

\begin{theorem}\label{thm:density}
Let $L\ge1$, $0<\rho<1$, $B\ge1$, and $0\le r\le B$ with
$\mathbb E_{u_L}r=1$. Set
\[
 e_\rho=\frac{1+\rho^2}{(1-\rho)^2},\qquad
 c(B,\rho)=B^{-1}(2B)^{-2e_\rho}.
\]
The two largest eigenvalues of $T_r$ satisfy
$1-\lambda_2/\lambda_1\ge c(B,\rho)>0$, independently of $L$.
In particular, no positive lower bound on $r$ is needed.
\end{theorem}
\begin{proof}
All norms and scalar products use $u_L$. The noise operator is positive
semidefinite, fixes constants and is a contraction. Hence
$\lambda_1\le B$. The unit vector $\sqrt r$ has Rayleigh quotient
$\langle r,P_\rho r\rangle\ge(\mathbb E r)^2=1$, by Walsh expansion,
so $1\le\lambda_1\le B$. On $S=\{r>0\}$ the matrix is strictly positive;
its Perron vector is positive on $S$ and its leading eigenvalue is simple.

We use two classical Boolean hypercontractive inequalities. Put
$q_*=1+\rho^{-2}$ and $p_*=1-\rho$. Then
\begin{equation}\label{eq:hyperinputs}
 \|P_\rho f\|_{q_*}\le\|f\|_2,\qquad
 \langle f,P_\rho g\rangle\ge\|f\|_{p_*}\|g\|_{p_*}\quad(f,g\ge0).
\end{equation}
For clarity about conventions, equations (1.1)--(1.2) and tensorization
in~\cite{MOS} use time $t=-\log\rho$ for the uniform Boolean semigroup.
For the second inequality set $p_*'=p_*/(p_*-1)=-(1-\rho)/\rho$.
Then
$\tfrac12\log((1-p_*')/(1-p_*))=-\log\rho$.
The reverse inequality gives $\|P_\rho g\|_{p_*'}\ge\|g\|_{p_*}$;
reverse H\"older gives the displayed bilinear bound. This is first applied
to positive functions, then to $f+\varepsilon,g+\varepsilon$ and the limit
$\varepsilon\downarrow0$. Finite sums and $p_*>0$ justify that limit;
we do not take a negative moment of a function with zeros.

If $\lambda_2<\lambda_1/2$, the gap exceeds $1/2$. Otherwise
$\lambda_2\ge1/2$. Choose a real unit second eigenvector $w$ orthogonal
to the Perron vector. It vanishes off $S$ and has both signs. Write
$w=s h_+-t h_-$, where $s,t>0$, $s^2+t^2=1$, and $h_+,h_-$ are
nonnegative orthonormal functions with disjoint supports. Positivity
and the eigenrelation give pointwise
$\lambda_2h_\pm\le T_rh_\pm$: on the corresponding sign support this
follows from the eigenrelation, and elsewhere from positivity.
Integration against the unit $h_\pm$ gives
$\lambda_2\le\langle\sqrt r h_\pm,P_\rho\sqrt r h_\pm\rangle$.
For $f_\pm=\sqrt r\,h_\pm$ this yields
\[
 \lambda_2f_\pm\le rP_\rho f_\pm,\qquad
 \|f_\pm\|_{q_*}\le\frac B{\lambda_2}\|f_\pm\|_2,\qquad
 \|f_\pm\|_2^2\ge\langle f_\pm,P_\rho f_\pm\rangle\ge\lambda_2.
\]
The middle inequality uses the first part of~\eqref{eq:hyperinputs}.
No division by $r$ occurs.

Let $\theta=p_*(q_*-2)/(2(q_*-p_*))$. Ordinary H\"older applied to
powers of $f$ gives
$\|f\|_2\le\|f\|_{p_*}^{\theta}\|f\|_{q_*}^{1-\theta}$,
since $1/2=\theta/p_*+(1-\theta)/q_*$. This does not require a normed
space at exponent $p_*<1$: the H\"older weights
$2\theta/p_*$ and $2(1-\theta)/q_*$ are positive and sum to one.
Thus each sign part satisfies
\[
 \|f_\pm\|_{p_*}\ge\sqrt{\lambda_2}
       (B/\lambda_2)^{-e_\rho},\qquad
 \frac{1-\theta}{\theta}
 =\frac{q_*(2-p_*)}{p_*(q_*-2)}=e_\rho.
\]
By the second part of~\eqref{eq:hyperinputs}, the cross entry obeys
\[
 d=\langle h_+,T_rh_-\rangle
   =\langle f_+,P_\rho f_-\rangle
   \ge\lambda_2(B/\lambda_2)^{-2e_\rho}.
\]
The compression of $T_r$ to the basis $(h_+,h_-)$ is
$\left(\begin{smallmatrix}a&d\\d&b\end{smallmatrix}\right)$.
Its eigenvector $(s,-t)$ has eigenvalue $\lambda_2$, giving
$a-\lambda_2=dt/s$ and $b-\lambda_2=ds/t$.
Its other eigenvalue is $\lambda_2+d/(st)$, because $s^2+t^2=1$.
The variational principle and $st\le1/2$ now imply
$\lambda_1-\lambda_2\ge d/(st)\ge2d$. Consequently
\[
 1-\frac{\lambda_2}{\lambda_1}
 \ge\frac{2\lambda_2}{\lambda_1}(B/\lambda_2)^{-2e_\rho}
 \ge B^{-1}(2B)^{-2e_\rho}.
\]
This does not deteriorate when one sign part has small norm. In the
previously omitted case, $c(B,\rho)\le1/4<1/2$, since $e_\rho\ge1$.
If the support has one point, $\lambda_2=0$ and the gap is one, so that
case is included as well.
\end{proof}

\begin{corollary}\label{cor:densitygrowth}
At fixed $\rho$, suppose a sequence has
$\gap(T_{r_L})\le K e^{-\gamma L}$, with $K,\gamma>0$. Then
\[
 \max r_L\ge
 \left(\frac{2^{-2e_\rho}}K\right)^{1/(2e_\rho+1)}
       e^{\gamma L/(2e_\rho+1)}.
\]
\end{corollary}
\begin{proof}
Apply Theorem~\ref{thm:density} with $B=\max r_L\ge1$ and rearrange.
\end{proof}
For Theorem~\ref{thm:main}, take $K=7$ and its positive $\gamma$.
The code mixture has maximum density $2^L(1-\eta)/|C|+\eta$, consistent
with this necessary growth. Upper density growth alone is not asserted
to be sufficient for collapse. Tiny minimum density alone cannot cause
collapse under a uniform maximum-density bound. These are statements
about this transfer family, not single-site heat-bath dynamics.
The constants are not claimed optimal.

The moment analysis separates universal guarantees from adverse examples.
The two strict exponent regimes are followed by auxiliary critical estimates
and an explicit rate, then the exact critical-budget theorem. The final
critical-tail section gives the sharp escaped-moment and fixed-cutoff
boundaries, and a stable
counterexample to necessity of uniform integrability. Earlier partial
estimates are retained for their explicit rate and proof dependencies;
they are not additional independent discoveries.

\section{Density-moment threshold: the two strict regimes}
\label{sec:finite-moment}
A pointwise density cap is stronger than necessary. We next replace it
by a dimension-independent finite moment. A second construction shows why
one cannot simply assume that any finite moment suffices. All norms in
this section use the uniform probability law $u_L$; the operator is the
same $T_r$ as in Section~\ref{sec:density}.

\begin{theorem}\label{thm:moment}
Fix $0<\rho<1$, $s>(1+\rho)/(1-\rho)$ and $B\ge1$. For every $L\ge1$
and density $r\ge0$ with $\mathbb E r=1$ and $\|r\|_s\le B$,
\[
 \gap(T_r)\ge\min\{1/2,(2B)^{-2e}\}>0,
 \qquad e=\frac{2\rho s}{(1-\rho)(s(1-\rho)-1-\rho)}.
\]
Neither a pointwise upper bound nor a positive lower bound is assumed.
\end{theorem}
\begin{proof}
Set $A=1+\rho$, $p=1-\rho$, $t=1+1/\rho$ and
\[
 Q=\frac{st}{s+t}=\frac{s(1+\rho)}{s\rho+1+\rho}.
\]
The condition on $s$ is exactly $Q>A$, so $Q>A>1>p>0$.
Ordinary Boolean hypercontractivity gives
\[
 \|P_{\sqrt\rho}f\|_2\le\|f\|_A,
 \qquad \|P_\rho f\|_t\le\|f\|_A.
\]
Indeed the noise thresholds are respectively
$\sqrt{(A-1)/(2-1)}=\sqrt\rho$ and
$\sqrt{(A-1)/(t-1)}=\rho$. The reverse bilinear inequality in
\eqref{eq:hyperinputs} applies at exponent $p$.

As before, testing the unit vector $\sqrt r$ gives $\lambda_1\ge1$.
For any real unit vector $h$, put $f=\sqrt r\,h$ and
$A_0=2s/(s+1)>A$. H\"older and probability-space monotonicity give
\[
 \langle h,T_rh\rangle=\|P_{\sqrt\rho}f\|_2^2
 \le\|f\|_A^2\le\|f\|_{A_0}^2
 \le\|\sqrt r\|_{2s}^2\|h\|_2^2\le B,
\]
where $1/A_0=1/(2s)+1/2$. Thus $1\le\lambda_1\le B$ without
using a maximum-density estimate.

If $\lambda_2<\lambda_1/2$, the gap is greater than $1/2$.
Otherwise $\lambda_2\ge1/2$. On the nonempty support of $r$ the matrix
has strictly positive entries. A real unit second eigenvector vanishes
off this support and has both signs, by orthogonality to the positive
Perron vector. Write it as $w=c h_+-d h_-$ with $c,d>0$,
$c^2+d^2=1$ and nonnegative orthonormal disjoint sign parts.
Positivity and the eigenrelation imply
$\lambda_2h_\pm\le T_rh_\pm$ pointwise, including their zero sets.
For either $f=\sqrt r\,h_\pm$, multiplication by $\sqrt r$ gives
$\lambda_2f\le rP_\rho f$. H\"older with $1/Q=1/s+1/t$ yields
\[
 \|f\|_Q\le (B/\lambda_2)\|f\|_A.
\]
Integration against the unit sign part also gives
$\lambda_2\le\langle f,P_\rho f\rangle\le\|f\|_A^2$.
Let $0<\theta<1$ solve $1/A=\theta/p+(1-\theta)/Q$.
H\"older applied to powers, with positive weights
$A\theta/p$ and $A(1-\theta)/Q$ summing to one, gives
$\|f\|_A\le\|f\|_p^\theta\|f\|_Q^{1-\theta}$.
No triangle inequality at exponent $p<1$ is used. Combining these bounds,
\[
 \|f\|_p\ge\sqrt{\lambda_2}(B/\lambda_2)^{-e},\qquad
 e=\frac{1-\theta}{\theta}
   =\frac{Q(A-p)}{p(Q-A)}
   =\frac{2\rho s}{(1-\rho)(s(1-\rho)-1-\rho)}.
\]
Reverse hypercontractivity now bounds the cross entry by
$D=\langle h_+,T_rh_-\rangle\ge\lambda_2(B/\lambda_2)^{-2e}$.
The two-dimensional compression, exactly as in the proof of
Theorem~\ref{thm:density}, has eigenvalues $\lambda_2$ and
$\lambda_2+D/(cd)$. Since $cd\le1/2$, the variational principle gives
\[
 \gap(T_r)\ge\frac{2D}{\lambda_1}
 \ge\frac{2\lambda_2}{\lambda_1}(B/\lambda_2)^{-2e}
 \ge(2B)^{-2e}.
\]
Here $2\lambda_2/\lambda_1\ge1$ and $\lambda_2\ge1/2$.
Combining the two cases proves the claim. One-point support has
$\lambda_2=0$ and is already covered. No step divides by $r$.
\end{proof}

\begin{theorem}\label{thm:tilts}
Fix $0<\rho<1$, $1<s<s_c:=(1+\rho)/(1-\rho)$ and $B>1$.
There are strictly positive mean-one densities $r_L$, for every $L\ge1$,
with $\|r_L\|_s\le B$ and $r_L\ge1/B$, and constants
$\kappa>1$ and $0<\omega<\kappa$, independent of $L$, such that
\[
 0<\gap(T_{r_L})\le
 2(\omega/\kappa)^L+\varepsilon^{-1}\kappa^{-L},
 \qquad \varepsilon=(B-1)/2.
\]
Thus a uniform positive lower density does not prevent exponential gap
collapse throughout the strict subcritical moment regime.
\end{theorem}
\begin{proof}
Relabel the cube as $\{-1,1\}^L$. For a fixed tilt $0<\delta<1$, put
\[
 n=\left(\frac{(1+\delta)^s+(1-\delta)^s}{2}\right)^{1/s},
 \qquad w_\pm(x)=\frac{1\pm\delta x}{n},\qquad m=n^{-1}.
\]
Strict convexity gives $n>1$, while $\|w_\pm\|_s=1$ and
$\mathbb E w_\pm=m$. Define $R_{\pm,L}(x)=\prod_i w_\pm(x_i)$ and
\[
 Z_L=1+2\varepsilon m^L,\qquad
 r_L=\frac{1+\varepsilon(R_{+,L}+R_{-,L})}{Z_L}.
\]
Product integration gives $\|R_{\pm,L}\|_s=1$ and
$\mathbb E R_{\pm,L}=m^L$. Hence $r_L>0$, $\mathbb E r_L=1$,
$\|r_L\|_s\le(1+2\varepsilon)/Z_L\le B$, and
$r_L\ge1/Z_L\ge1/B$.

Write $Q=P_{\sqrt\rho}$. Since $Q^2=P_\rho$, the matrices
$T_{r_L}=AA^*$ and $S_L=A^*A$, with $A=M_{\sqrt{r_L}}Q$,
have identical spectra, including multiplicities. Here $M_f$ denotes
multiplication by $f$. Direct expansion gives
\[
 Z_L S_L=P_\rho+\varepsilon(X_L+Y_L),\qquad
 X_L=H_+^{\otimes L},\quad Y_L=H_-^{\otimes L},
\]
where the one-site matrices, in the orthonormal Walsh basis $(1,x)$, are
\[
 H_\pm=\frac1n
 \begin{pmatrix}1&\pm\delta\sqrt\rho\\
                 \pm\delta\sqrt\rho&\rho\end{pmatrix}.
\]
They are positive definite and conjugate by $\diag(1,-1)$. Their common
simple top eigenvalue is
\[
 \kappa(\delta)=
 \frac{1+\rho+\sqrt{(1-\rho)^2+4\rho\delta^2}}{2n(\delta)}.
\]
At zero, $\kappa(0)=1$, $(\log\kappa)'(0)=0$, and differentiation yields
\[
 (\log\kappa)''(0)=\frac{2\rho}{1-\rho}-(s-1)>0.
\]
All functions are smooth near zero. Taylor's theorem therefore gives
$\kappa>1$ for every sufficiently small positive $\delta$. Fix one such
$\delta<1$. This choice depends on $\rho,s$ and is independent of $L$;
no uniform exponential rate as $s\uparrow s_c$ is asserted.

Set $\omega=\|H_+^{1/2}H_-^{1/2}\|$. Positive definiteness and
submultiplicativity give $0<\omega\le\kappa$. If equality held, norm
attainment at a unit vector $v$ would force equality in both steps of
\[
 \|H_+^{1/2}H_-^{1/2}v\|
 \le\sqrt\kappa\,\|H_-^{1/2}v\|\le\kappa.
\]
The second equality makes $v$ a top eigenvector of $H_-$; the first
then makes it a top eigenvector of $H_+$. But $H_+-H_-$ has zero
diagonal and nonzero equal off-diagonal entries, so is invertible.
It cannot annihilate their common top eigenvector. Thus $\omega<\kappa$.

For $N=2^L\ge2$, use the $N$-by-$2N$ matrix
$C=[\sqrt{X_L},\sqrt{Y_L}]$. Its two Gram matrices satisfy
\[
 CC^*=X_L+Y_L,\qquad
 C^*C=\begin{pmatrix}
 X_L&\sqrt{X_L}\sqrt{Y_L}\\
 \sqrt{Y_L}\sqrt{X_L}&Y_L
 \end{pmatrix}.
\]
The block diagonal part has its two largest eigenvalues equal to
$\kappa^L$. Tensorization of positive square roots and of operator
norms gives $\|\sqrt{X_L}\sqrt{Y_L}\|=\omega^L$, which is also the
norm of the self-adjoint off-diagonal perturbation. The variational
eigenvalue bound consequently gives
\[
 \lambda_1(X_L+Y_L)\le\kappa^L+\omega^L,\qquad
 \lambda_2(X_L+Y_L)\ge\kappa^L-\omega^L>0.
\]
Indeed $C^*C$ and $CC^*$ have the same $N$ positive eigenvalues,
with $N$ additional zeros only in $C^*C$, so both leading bounds
transfer even when $L=1$. Also $X_L+Y_L\ge X_L$ gives
$\lambda_1(X_L+Y_L)\ge\kappa^L$.

Finally $P_\rho$ is positive semidefinite of norm one. For
$F_L=Z_L S_L$, the variational principle and norm bound imply
\[
 \lambda_1(F_L)\le\varepsilon(\kappa^L+\omega^L)+1,
 \qquad \lambda_2(F_L)\ge\varepsilon(\kappa^L-\omega^L),
 \qquad \lambda_1(F_L)\ge\varepsilon\kappa^L.
\]
Subtracting and dividing gives the claimed upper bound. The scalar
$Z_L$ cancels from the ratio, and $S_L$ and the actual $T_{r_L}$
have the same spectrum. Its entrywise positivity gives strict Perron
separation at every finite $L$. Both terms in the upper bound decay
exponentially, although the bound can exceed one at small $L$.
\end{proof}

Theorems~\ref{thm:moment} and~\ref{thm:tilts} separate the two strict
regimes at the same value $s_c$. The next two sections treat the critical exponent $s=s_c$:
first quantitative bounds in two budget ranges, then the exact budget
boundary including equality. At $B=1$, Jensen forces $r=1$ and the actual gap is $1-\rho$;
the sufficient lower bound is not claimed optimal. The tilted family
does not claim product marginals for its weighting or stationary law.
It isolates the moment obstruction for the same transfer operator,
whereas Theorem~\ref{thm:main} isolates the marginal-information
obstruction. The earlier two-point-spike construction in the retained
analytical sources covers the narrower interval
$1<s<\log2/\log(2/(1+\rho))$; it remains valid but is superseded here
by the product-tilt construction.


\section{Auxiliary critical estimates and an explicit rate}
At $s=s_c=(1+\rho)/(1-\rho)$ the size of the moment budget matters.
We give a sufficient stability condition near $B=1$ and a counterexample
for every $B>2$. These preliminary bounds are retained for their quantitative content.
Section~\ref{sec:critical-exact} then resolves the intervening budget
range and the exact boundary. All norms in this section use the uniform
probability measure $u_L$.

\begin{theorem}[Small critical budgets]\label{thm:critical-small}
Fix $0<\rho<1$ and $s=s_c$. For any $L\ge1$ and any density $r\ge0$
with $\mathbb E r=1$ and $\|r\|_s\le B$, $B\ge1$,
\[
 \gap(T_r)\ge\max\{0,1-\rho-d_s(B)\},\qquad
 d_s(B)=\begin{cases}
 (B^s-1)^{1/s},&s\ge2,\\
 \sqrt{\dfrac{2(B^s-1)}{s(s-1)}}(B+1)^{(2-s)/2},&1<s<2.
 \end{cases}
\]
In particular, each fixed $\rho$ admits a nontrivial interval of budgets
$B>1$ with a positive dimension-independent gap. No pointwise density
bounds are required.
\end{theorem}
\begin{proof}
Put $Q=P_{\sqrt\rho}$ and $S_r=QM_rQ$, which has the spectrum of $T_r$.
Ordinary Boolean hypercontractivity, recalled in~\cite[Eq.~(1.1)]{MOS},
gives $\|Qf\|_q\le\|f\|_2$ for $q=2s/(s-1)$, since
$\rho=(s-1)/(s+1)=1/(q-1)$. H\"older's inequality with exponents
$s,q,q$ yields
\[
 |\langle f,Q M_{r-1}Qg\rangle|
 \le\|r-1\|_s\|Qf\|_q\|Qg\|_q
 \le\|r-1\|_s\|f\|_2\|g\|_2.
\]
Thus $\|S_r-P_\rho\|\le\|r-1\|_s$ and
$\lambda_2(S_r)\le\rho+\|r-1\|_s$. Also
$\lambda_1(S_r)\ge\langle1,S_r1\rangle=1$.
It remains to show $\|r-1\|_s\le d_s(B)$.

Define $\phi_s(t)=t^s-1-s(t-1)$ for $t\ge0$, so that
$\mathbb E\phi_s(r)\le B^s-1$. For $s\ge2$ and $u\ge0$,
\[
 \phi_s(1+u)=s\int_0^u((1+v)^{s-1}-1)\,dv\ge u^s,
\]
using $(1+v)^{s-1}\ge1+v^{s-1}$. For $0\le u\le1$,
the derivative of $\phi_s(1-u)-u^s$ is
$s[1-(1-u)^{s-1}-u^{s-1}]\ge0$ and its value at zero is zero.
Hence $\phi_s(t)\ge|t-1|^s$, proving the first estimate by integration.
For $1<s<2$, the integral Taylor remainder gives
\begin{align*}
 \phi_s(t)&=s(s-1)(t-1)^2
  \int_0^1(1-v)[1+v(t-1)]^{s-2}\,dv\\
 &\ge\frac{s(s-1)}2\frac{(t-1)^2}{(t+1)^{2-s}}.
\end{align*}
At $t=0$ the identity follows by continuity; the integral is finite for
$s>1$. For $W=(r-1)^2/(r+1)^{2-s}$, H\"older gives
\[
 \mathbb E|r-1|^s
 \le(\mathbb EW)^{s/2}
       (\mathbb E(r+1)^s)^{(2-s)/2}.
\]
Use $\mathbb EW\le2(B^s-1)/(s(s-1))$ and $\|r+1\|_s\le B+1$.
Both formulas for $d_s(B)$ tend to zero as $B\downarrow1$.
Positive semidefiniteness supplies the trivial zero lower bound.
\end{proof}

\begin{theorem}[Balanced critical tilt]\label{thm:critical-large}
Fix $0<\rho<1$, $s=s_c$, and $B>2$. There are constants $c,d>0$,
independent of $L$, and strictly positive densities $r_L$, for every $L\ge1$,
such that
\[
 \mathbb E r_L=1,\quad\|r_L\|_s\le B,\quad r_L\ge1/B,
 \qquad
 0<\gap(T_{r_L})\le2\sqrt{B-1}\,e^{-cd^2L^{1/3}}.
\]
\end{theorem}
\begin{proof}
For a uniform sign $x\in\{-1,1\}$ and $|\delta|<1$ define
\[
 n(\delta)=\left(\frac{(1+\delta)^s+(1-\delta)^s}{2}\right)^{1/s},
 \quad w_\delta(x)=\frac{1+\delta x}{n(\delta)},\quad
 m(\delta)=n(\delta)^{-1}.
\]
The one-site matrix is
\[
 H(\delta)=\frac1{n(\delta)}
 \begin{pmatrix}1&\delta\sqrt\rho\\\delta\sqrt\rho&\rho\end{pmatrix},
 \qquad D=\diag(1,\rho).
\]
It is positive definite with top eigenvalue
\[
 \kappa(\delta)=
 \frac{1+\rho+\sqrt{(1-\rho)^2+4\rho\delta^2}}{2n(\delta)}.
\]
For $\tau(\delta)=\|D^{1/2}H(\delta)^{1/2}\|$, the square is the
top eigenvalue of $D^{1/2}H(\delta)D^{1/2}$, so
\[
 \tau(\delta)^2=
 \frac{1+\rho^2+\sqrt{(1-\rho^2)^2+4\rho^2\delta^2}}{2n(\delta)}.
\]
These positive functions have even analytic logarithms near zero.
Taylor expansion with bounded fourth derivative gives, at $s=s_c$,
\begin{align*}
 \log n(\delta)&=\frac{\rho}{1-\rho}\delta^2+O(\delta^4),\\
 \log\kappa(\delta)&=O(\delta^4),\\
 \log\tau(\delta)&=-\frac{\rho}{2(1-\rho^2)}\delta^2+O(\delta^4).
\end{align*}
Indeed, the respective square-root numerators divided by two are
$1+\rho\delta^2/(1-\rho)+O(\delta^4)$ and
$1+\rho^2\delta^2/(1-\rho^2)+O(\delta^4)$.
Consequently there are $0<\delta_*<1$, $C\ge1$ and $c>0$ such that,
for $|\delta|\le\delta_*$,
\[
 |\log\kappa(\delta)|\le C\delta^4,
 \qquad\log\tau(\delta)\le-c\delta^2.
\]
For example, $c=\rho/[4(1-\rho^2)]$ works after reducing $\delta_*$.
Choose a fixed $0<d\le\delta_*$ with $Cd^4\le\log(B-1)$, possible
because $B>2$. Put
\[
 \delta_L=dL^{-1/3},\quad a_L=\kappa(\delta_L)^{-L},\quad
 R_L=\prod_{i=1}^Lw_{\delta_L}(x_i),\quad
 Z_L=1+a_Lm(\delta_L)^L,\quad r_L=\frac{1+a_LR_L}{Z_L}.
\]
The coefficient balances the two top block eigenvalues exactly.
Since $|\log a_L|\le Cd^4L^{-1/3}\le\log(B-1)$,
we have $a_L\le B-1$ and $a_L\to1$.
Product measure gives $\|R_L\|_s=1$ and
$\mathbb ER_L=m(\delta_L)^L$, with $0<m(\delta_L)<1$.
Therefore $\mathbb Er_L=1$, $\|r_L\|_s\le(1+a_L)/Z_L\le B$,
and $r_L\ge1/Z_L\ge1/B$ for every $L\ge1$.

Let $X_L=H(\delta_L)^{\otimes L}$. Spectral equivalence gives
$F_L=Z_LS_{r_L}=P_\rho+a_LX_L$. Both summands have top eigenvalue one.
For the $N$ by $2N$ matrix
$\mathcal C_L=[P_\rho^{1/2},\sqrt{a_L}X_L^{1/2}]$, $N=2^L$,
the Gram matrix $\mathcal C_L^*\mathcal C_L$ has those summands on
its diagonal and off-diagonal block of norm
\[
 \eta_L=\sqrt{a_L}\,\tau(\delta_L)^L
 \le\sqrt{B-1}\,e^{-cd^2L^{1/3}}.
\]
Its block diagonal has two leading eigenvalues exactly one. The
off-diagonal self-adjoint perturbation has norm $\eta_L$, whence
$\lambda_1\le1+\eta_L$ and $\lambda_2\ge1-\eta_L$.
The $N$ positive eigenvalues of $\mathcal C_L\mathcal C_L^*=F_L$
are precisely the positive eigenvalues of $\mathcal C_L^*\mathcal C_L$;
the latter has $N$ additional zeros. Thus the bounds apply to the first
two eigenvalues of $F_L$, including at $L=1$. Since $F_L\ge P_\rho$,
$\lambda_1(F_L)\ge1$ and $\gap(F_L)\le2\eta_L$.
Scalar normalization and $AA^*/A^*A$ preserve the ratio. Strict entrywise
positivity of $T_{r_L}$ gives its strictly positive finite-dimensional gap.
\end{proof}

The uniform component is a spectral channel here, rather than an additive
error to discard. At the critical exponent the tilted top eigenvalue has
zero quadratic variation, while its cross norm with the uniform channel
has negative quadratic variation. The dimension-dependent tilt separates
these orders, and balancing removes the residual diagonal mismatch.
The small-budget estimate uses classical hypercontractivity and elementary
perturbation theory; the large-budget family asserts neither matching
marginals nor fixed-range interactions. The sufficient
budget $B>2$ is not optimal; the next section identifies the exact
threshold without claiming a rate there.

\section{The exact critical moment budget}\label{sec:critical-exact}
The preceding balanced construction has a sharper limiting moment than
the triangle inequality used to obtain its all-dimension rate bound.
Combining that limit with a sign decomposition of a second eigenvector
identifies the entire critical budget boundary. Throughout this section,
all norms and expectations use uniform probability on $\{-1,1\}^L$.

\begin{theorem}[Exact critical budget]\label{thm:critical-exact}
Fix $0<\rho<1$, $s=(1+\rho)/(1-\rho)$, and $B_*=2^{1/s}$.
For every $L\ge1$ and $r\ge0$ with $\mathbb Er=1$,
\begin{equation}\label{eq:two-sign-budget}
 2\lambda_2(T_r)^s\le\mathbb Er^s.
\end{equation}
Consequently $\|r\|_s\le B<B_*$ implies
\[
 \gap(T_r)\ge 1-B/B_*>0.
\]
For each fixed $b\in(0,1)$ there are strictly positive densities
$\widetilde r_L$, for every $L\ge1$, such that
\[
 \mathbb E\widetilde r_L=1,\qquad
 \|\widetilde r_L\|_s\le B_*,\qquad
 \widetilde r_L\ge b,\qquad
 \gap(T_{\widetilde r_L})\longrightarrow0.
\]
Thus the infimum of the gap over all dimensions and densities with
$\|r\|_s\le B$ is positive exactly for $1\le B<B_*$.
The same classification holds with any prescribed fixed floor
$r\ge b$, $0<b<1$. The positive lower bound displayed above is not
asserted optimal; no convergence rate at $B=B_*$ is asserted.
\end{theorem}

\begin{proof}
\emph{Two sign regions.}
Put $Q=P_{\sqrt\rho}$ and $p=1+\rho=2s/(s+1)$.
Endpoint Boolean hypercontractivity gives $\|Qg\|_2\le\|g\|_p$.
For $h$ supported on $C$, H\"older's inequality gives
\begin{align}
 \langle h,T_rh\rangle
 &=\|Q(\sqrt r\,h)\|_2^2\nonumber\\
 &\le\|\sqrt r\,h\|_p^2
 \le\bigl(\mathbb E[r^s\mathbf1_C]\bigr)^{1/s}\|h\|_2^2.
 \label{eq:restricted-moment}
\end{align}
Indeed, before raising to $2/p$, the conjugate exponents are
$2/(2-p)$ and $2/p$, and $p/(2-p)=s$.

On its nonempty active support $A=\{r>0\}$, the matrix $T_r$
is strictly entrywise positive. It has a simple positive Perron root
and a Perron vector strictly positive on $A$. Its other coordinates
form a zero block. If $\lambda_2=0$,~\eqref{eq:two-sign-budget}
is immediate. Otherwise a real second eigenvector $f$ vanishes off
$A$ and is orthogonal to the Perron vector. Both its positive and
negative parts $f_+,f_-$ are therefore nonzero, including when
$\lambda_2$ has multiplicity greater than one.
Entrywise positivity gives
\[
 \lambda_2\|f_+\|_2^2
 =\langle f_+,T_rf_+\rangle-\langle f_+,T_rf_-\rangle
 \le\langle f_+,T_rf_+\rangle.
\]
Applying~\eqref{eq:restricted-moment} on $\{f>0\}$ yields
$\lambda_2^s\le\mathbb E[r^s\mathbf1_{\{f>0\}}]$.
The negative part gives the analogous bound on $\{f<0\}$.
Adding over these disjoint regions proves~\eqref{eq:two-sign-budget}.
No mean-zero condition under uniform measure is imposed on $f$.
Finally $\sqrt r$ is a unit vector and
\[
 \lambda_1(T_r)\ge\langle r,P_\rho r\rangle
 \ge(\mathbb Er)^2=1
\]
by the nonnegative Walsh eigenvalues. This proves the positive gap
bound below $B_*$. At $B=1$, strict Jensen forces $r=1$ and the
actual gap is $1-\rho$.

\emph{The balanced family's exact moment.}
Use the single-site quantities $n,w,m,H,\kappa,\tau$ from
Theorem~\ref{thm:critical-large}, with any sufficiently small fixed
$d>0$ and $\delta_L=dL^{-1/3}$. This step does not impose the
earlier condition involving $B>2$. Define
\[
 a_L=\kappa(\delta_L)^{-L},\quad
 R_L=\prod_{i=1}^Lw_{\delta_L}(x_i),\quad
 Z_L=1+a_Lm(\delta_L)^L,\quad
 r_L=(1+a_LR_L)/Z_L.
\]
The critical expansions already proved give $a_L\to1$ and
$m(\delta_L)^L\to0$, hence $Z_L\to1$. For the latter limit,
$\log m(\delta)=-(s-1)\delta^2/2+O(\delta^4)$ and
$L\delta_L^2\to\infty$, whereas $L\delta_L^4\to0$.
The overlap $\eta_L=\sqrt{a_L}\tau(\delta_L)^L$ tends to zero.
More explicitly, the quadratic coefficient of $\log\tau$ is
\[
 \frac12\left[\frac{\rho^2}{1-\rho^2}-\frac{s-1}{2}\right]
 =-\frac{\rho}{2(1-\rho^2)}.
\]
The same Gram comparison, now without a prescribed all-$L$ budget,
gives $1\le\lambda_1(S_L)\le1+\eta_L$ and
$\lambda_2(S_L)\ge1-\eta_L$ for
$S_L=P_\rho+a_LH(\delta_L)^{\otimes L}$.
For large $L$, the second bound is positive, so it corresponds to
a nonzero eigenvalue rather than an added Gram zero.
Since $QM_{r_L}Q=S_L/Z_L$, both of its top eigenvalues tend to one.
They are the actual eigenvalues of $T_{r_L}$: for
$X=M_{\sqrt{r_L}}Q$, the square matrices $X^*X$ and $XX^*$
have identical spectra, including zeros and multiplicities.

For each fixed $0<t<s$, direct single-site expansion yields
\[
 \log\mathbb Ew_\delta^t
 =\frac{t(t-s)}2\delta^2+O(\delta^4).
\]
Thus $\mathbb ER_L^t\to0$, while $\mathbb ER_L^s=1$ exactly.
The remainders may depend on fixed $t$; no uniformity near the
endpoints is needed. Set $\theta=\min\{1,s/2\}$.
Both $\theta$ and $s-\theta$ belong to $(0,s)$, and
\begin{equation}\label{eq:moment-cross}
 0\le(1+y)^s-1-y^s
 \le s2^{s-1}\bigl(y^\theta+y^{s-\theta}\bigr),\qquad y\ge0.
\end{equation}
For $0\le y\le1$, the mean-value theorem bounds the left side by
$s2^{s-1}y\le s2^{s-1}y^\theta$. For $y\ge1$, apply that
bound to $1/y$ and multiply by $y^s$, obtaining
$s2^{s-1}y^{s-1}\le s2^{s-1}y^{s-\theta}$.
The lower bound follows from $s>1$.
Substitute $y=a_LR_L$ in~\eqref{eq:moment-cross} and take expectations.
The two interior moments vanish and $a_L\to1$, so
\[
 \mathbb E(1+a_LR_L)^s=1+a_L^s+o(1)\longrightarrow2.
\]
Dividing by $Z_L^s\to1$ proves
$B_L:=\|r_L\|_s\to2^{1/s}=B_*$.

\emph{The exact boundary and a prescribed floor.}
Each $r_L$ is nonconstant, so $B_L>1$. Put
\[
 \varepsilon_L=\begin{cases}
 0,&B_L\le B_*,\\
 (B_L-B_*)/(B_L-1),&B_L>B_*,
 \end{cases}
 \qquad \widetilde r_L=(1-\varepsilon_L)r_L+\varepsilon_L.
\]
The only denominator used is positive. Since $B_L\to B_*>1$,
we have $0\le\varepsilon_L<1$ and $\varepsilon_L\to0$.
Normalization and positivity persist, and Minkowski gives
\[
 \|\widetilde r_L\|_s
 \le(1-\varepsilon_L)B_L+\varepsilon_L\le B_*.
\]
Because $Z_L>1$, both $r_L$ and $1$ are at least $1/Z_L$,
so the mixture has that same floor.

For signed $h$, the dual endpoint estimate
$\|Qf\|_q\le\|f\|_2$, $q=2s/(s-1)$, and H\"older with
$1/s+2/q=1$ give $\|QM_hQ\|\le\|h\|_s$. Hence
\[
 \|QM_{\widetilde r_L-r_L}Q\|
 \le\varepsilon_L\|1-r_L\|_s
 \le\varepsilon_L(1+B_L)\longrightarrow0.
\]
Weyl's bound and the exact $X^*X/XX^*$ spectral identity imply
that both top eigenvalues of $T_{\widetilde r_L}$ still tend to one.
Its normalized gap therefore tends to zero. This argument requires
neither pointwise upper control nor a square-root multiplier estimate.

Given fixed $0<b<1$, eventually $1/Z_L\ge b$. Replace the finite
initial segment by the constant density $1$. Every stated constraint
then holds for every $L\ge1$, with the same limiting gap. The same
sequence is admissible for all $B\ge B_*$. Each finite transfer
matrix remains strictly positive, so its individual gap is positive;
only the dimension-free infimum vanishes. A floor $b=1$ would force
the uniform density and is explicitly excluded.
\end{proof}

The sign-region argument is classical. In the additive lattice setting,
Biskup and K\"onig~\cite[Lemma~4.5 and Eq.~(4.29)]{BiskupKoenig}
bound the second eigenvalue using the two sign regions and split a
potential budget between them. Here endpoint hypercontractivity gives
the restricted $s$-moment estimate, while the balanced construction
and vanishing dilution attain the boundary asymptotically. This
comparison credits the mechanism; it does not establish historical
priority of the exact multiplicative-cube classification. The new
family makes no matching-marginal or fixed-range interaction claim.


% UNCOMPILED DRAFT for a grouped revision; requires source302 audit and document review.
\section{The escaped critical-moment boundary}
\label{sec:critical-tails}
Fix $0<\rho<1$ and $s=(1+\rho)/(1-\rho)$ throughout this section.
Expectations and norms refer to the uniform Boolean measure. Write
$Q=P_{\sqrt\rho}$ and $T_r=M_{\sqrt r}P_\rho M_{\sqrt r}$.
We also use $\gap(r):=\gap(T_r)$ when the density shorthand is convenient.
Here
$r\ge0$ and $\mathbb E r=1$. For a family indexed by $L\ge1$, define
\[
 H(K)=\sup_L\mathbb E[r_L^s\mathbf1_{r_L>K}],\qquad
 \Theta=\lim_{K\to\infty}H(K)=\inf_{K\ge1}H(K).
\]
The nonnegative monotone limit is allowed to be infinite. Uniform
integrability of the critical powers is the special case $\Theta=0$.

\begin{theorem}[Finite cutoff with tail budget below one]\label{thm:critical-tail}
Suppose $K\ge1$, $0\le\tau<1$, and
$\mathbb E[r^s\mathbf1_{r>K}]\le\tau$. Put
\[
 t=\tau^{1/(2s)},\quad m=\frac{1+t}{2},\quad d=\frac{1-t}{2},\quad
 B_K=(K^{s-1}+\tau)^{1/s}.
\]
Then
\[
 \gap(T_r)\ge
 \min\left\{1-m^2,\frac{2d^{2s^2}K^{1-s^2}}{B_K}\right\}>0.
\]
In particular, $\Theta<1$ implies a dimension-independent positive gap.
No optimality is claimed for the displayed positive constant.
\end{theorem}
\begin{proof}
For $w\ge0$ and $\|h\|_2=1$, H\"older and forward hypercontractivity,
with $p=1+\rho$ and $p/(2-p)=s$, give
\[
 \|Q\sqrt w h\|_2^2\le\|\sqrt w h\|_p^2\le\|w\|_s.
\]
Since $r^s\mathbf1_{r\le K}\le K^{s-1}r$, the total moment is at most
$K^{s-1}+\tau$. Thus $\lambda_1(T_r)\le B_K$. The unit vector
$\sqrt r$ also gives $\lambda_1\ge\langle r,P_\rho r\rangle\ge1$
by the nonnegative Walsh spectrum.
If the active support has one point, $\lambda_2=0$ and the gap is one.
If $\lambda_2<m^2$, the gap is greater than $1-m^2$.

Otherwise choose a real second eigenvector, orthogonal to the positive
active Perron vector. Write it as $u h_+-v h_-$, with $u,v>0$,
$u^2+v^2=1$, and disjointly supported nonnegative unit vectors
$h_+,h_-$. Their uniform means need not be equal. Set
$D=\langle h_+,T_rh_-\rangle>0$. The eigenvector equations give
\[
 A_+=\lambda_2+(v/u)D,\qquad A_-=\lambda_2+(u/v)D,
 \quad A_\pm=\langle h_\pm,T_rh_\pm\rangle.
\]
For each sign split $f=\sqrt r h_\pm$ at $r=K$. Then
$\|Qf\|_2\ge m$ and $\|Qf_{\rm high}\|_2\le t$, so
\[
 \|Qf_{\rm low}\|_2\ge d,\qquad
 \|f_{\rm low}\|_p\ge d,\qquad
 \|f_{\rm low}\|_2\le\sqrt K.
\]
For $a=1-\rho$ and $\theta=s^{-2}$ the identity
$1/p=\theta/a+(1-\theta)/2$ and power-H\"older give
$\|g\|_p\le\|g\|_a^\theta\|g\|_2^{1-\theta}$.
This applies ordinary H\"older to powers and requires no quasi-norm
triangle inequality. The positive lower $p$-norm permits division;
hence $\|f_{\rm low}\|_a\ge d^{1/\theta}K^{-(1-\theta)/(2\theta)}$.
Positivity and reverse hypercontractivity yield
\[
 D\ge\langle f_{+,\rm low},P_\rho f_{-,\rm low}\rangle
 \ge d^{2s^2}K^{1-s^2}.
\]
The sign compression has eigenvalues $\lambda_2$ and
$\lambda_2+D/(uv)$. Its largest eigenvalue is at most $\lambda_1$,
so $\lambda_1-\lambda_2\ge2D$. Dividing by $\lambda_1\le B_K$
proves the bound, including repeated second eigenvalues.
Finally, if $\Theta<1$, choose $\Theta<\tau<1$ and a single finite
$K$ with $H(K)\le\tau$. The same constants then apply to every $L$.
\end{proof}

Uniform integrability is sufficient but stronger than needed. For example,
a bound
\[
 \sup_L\mathbb E[r_L^s\psi(r_L)]\le C<\infty
\]
with a
nonnegative nondecreasing function $\psi$ diverging at infinity implies
$H(K)\le C/\psi(K)$ when $\psi(K)>0$. This is a weighted-moment
condition; an arbitrary such $\psi$ need not be a Young function or
define a Luxemburg norm. A fixed $s+\varepsilon$ moment bound also
implies $H(K)\le K^{-\varepsilon}\sup_L\mathbb E r_L^{s+\varepsilon}$.
Neither condition follows from a mere critical-moment bound.

A total bound $\sup_L\|r_L\|_s\le B$ does imply
\begin{equation}\label{eq:tail-total-budget}
 \Theta\le B^s-1.
\end{equation}
Indeed $\mathbb E[r_L\mathbf1_{r_L>K}]\le B^sK^{1-s}$, while Jensen
for $X_L=r_L\mathbf1_{r_L\le K}$ gives
$\mathbb E X_L^s\ge(1-B^sK^{1-s})_+^s$. Subtraction and the limit
$K\to\infty$ prove~\eqref{eq:tail-total-budget}. For $B<2^{1/s}$ this
recovers a positive gap. The direct exact-budget theorem already gives
the stronger explicit lower bound $1-B/2^{1/s}$.


\begin{proposition}[Escaped moment with a stable gap]\label{prop:stable-spike}
For any $M>1$, $B>M^{1/s}$, and $0<b<1$, there are densities $r_L$
with $\mathbb E r_L=1$, $r_L\ge b$, $\|r_L\|_s\le B$, and
\[
 \mathbb E r_L^s\to M,\qquad
 \mathbb E[r_L^s\mathbf1_{r_L>K}]\to M-1\quad(K>1\text{ fixed}),
\]
while $\gap(T_{r_L})\to1-\rho$ and
$\gap(T_{r_L})\ge(1-\rho)/2$ for every $L\ge1$.
\end{proposition}
\begin{proof}
Set $N=2^L$, $c=(M-1)^{1/s}$, choose a vertex $z_L$, and first put
\[
 q_L=\frac{1+cN^{1/s}\mathbf1_{z_L}}{Z_L},\qquad
 Z_L=1+cN^{1/s-1}.
\]
These densities have mean one, lower bound $1/Z_L\to1$, and moments
\[
 \mathbb E q_L^s=
 \frac{1-N^{-1}+(N^{-1/s}+c)^s}{Z_L^s}\longrightarrow1+c^s=M.
\]
For each fixed $K>1$ only the exceptional vertex eventually exceeds
$K$, and its contribution tends to $c^s=M-1$.

The operators $T_{q_L}$ and $S_L=QM_{q_L}Q$ are isospectral, including
multiplicities, since they are $AA^*$ and $A^*A$ for
$A=M_{\sqrt{q_L}}Q$. In probability-normalized $L^2$, the vector
$e_z=\sqrt N\mathbf1_{z_L}$ is a unit vector and
$M_{\mathbf1_{z_L}}=|e_z\rangle\langle e_z|$. Therefore
\[
 S_L=\frac{P_\rho+cN^{1/s}Q|e_z\rangle\langle e_z|Q}{Z_L},
 \qquad
 \|Q|e_z\rangle\langle e_z|Q\|
 =\langle e_z,P_\rho e_z\rangle
 =\left(\frac{1+\rho}{2}\right)^L.
\]
It follows that $\|S_L-P_\rho\|\le\delta_L$, where
\[
 \delta_L=\frac{c d_\rho^L+c a_\rho^L}{Z_L},\quad
 d_\rho=2^{-(1-1/s)},\quad
 a_\rho=\frac{1+\rho}{2^{2\rho/(1+\rho)}}.
\]
Both bases lie strictly between zero and one. For $a_\rho$, strict
convexity of $t\log t$ on $[1,2]$ gives
$(1+\rho)\log(1+\rho)<2\rho\log2$.

The largest two eigenvalues of $P_\rho$ are $1$ and $\rho$, counting
multiplicity. The min-max principle bounds the changes in both by
$\delta_L$. Also $\langle1,S_L1\rangle=1$, so $\lambda_1(S_L)\ge1$.
Consequently
\[
 \left|\frac{\lambda_2(S_L)}{\lambda_1(S_L)}-\rho\right|
 \le(1+\rho)\delta_L.
\]
Choose $L_0$ so that for every $L\ge L_0$ the floor and norm bounds
hold and $c(d_\rho^L+a_\rho^L)\le(1-\rho)/(2(1+\rho))$. Existence
uses $B^s>M$ and the strict exponential contractions. Finally set
$r_L=1$ for $L<L_0$ and $r_L=q_L$ otherwise. The prefix has gap
$1-\rho$; the tail has the asserted common bound. All limits persist.
\end{proof}

\begin{proposition}[Sharpness and nonnecessity of the tail criterion]
\label{prop:sharp-tail}
The universal guarantee $\Theta<1$ cannot be extended to $\Theta\le1$:
the exact-budget family of Theorem~\ref{thm:critical-exact} has
$\Theta=1$ and gaps tending to zero. For every finite $\Theta_0>0$,
there are nevertheless families with $\Theta=\Theta_0$ and a uniform
positive gap, even with an arbitrarily prescribed density floor below
one. Thus $\Theta\ge1$ is necessary for collapse, but not sufficient.
\end{proposition}
\begin{proof}
First suppose $r_L\to1$ in probability and $\mathbb E r_L^s\to M<\infty$.
For each fixed $K>1$, the function $F_K(x)=x^s\mathbf1_{x\le K}$ is
bounded and continuous at one, so $\mathbb E F_K(r_L)\to1$.
This follows by splitting a neighborhood of one from its probability-small
complement and works on the varying cubes. Therefore the tail moments
tend to $M-1$ for every such $K$, giving $\Theta\ge M-1$.
For the upper bound fix $\varepsilon>0$ and $K_0>1$. For all sufficiently
large $L$, the total moment is at most $M+\varepsilon$ and
$\mathbb E F_{K_0}(r_L)\ge1-\varepsilon$. For every $K\ge K_0$ their
tail moments are then at most $M-1+2\varepsilon$. A larger cutoff kills
the tails of the finitely many remaining finite functions. Taking the
supremum over all $L$, then $K\to\infty$ and $\varepsilon\downarrow0$,
proves $\Theta=M-1$. This finite-prefix step is required; the fixed-$K$
limit alone would not justify exchanging the supremum and cutoff limit.

The diluted balanced family of Theorem~\ref{thm:critical-exact} satisfies
$r_L\to1$ in probability: its tilt has every smaller positive moment
tending to zero, while the balancing and dilution parameters tend to one
and zero, respectively. Also $\lambda_2(T_{r_L})\to1$ and
$2\lambda_2^s\le\mathbb E r_L^s\le2$, so the critical moment tends to
two. Hence $\Theta=1$, while that family's gap tends to zero.
Finally apply Proposition~\ref{prop:stable-spike} with
$M=1+\Theta_0$. Its densities converge to one in probability, have
critical moment tending to $M$, and keep the positive gap there proved.
Thus $\Theta=\Theta_0$. Its strict total-budget slack and finite-prefix
replacement remain part of the construction.
\end{proof}

\begin{theorem}[The fixed-cutoff endpoint]\label{thm:finite-tail-endpoint}
For every $0<b<1$, there is a single family of densities $r_L$ with
\[
 \mathbb E r_L=1,\qquad r_L\ge b,\qquad \mathbb E r_L^s\le2,
 \qquad \mathbb E[r_L^s\mathbf1_{r_L>1}]\le1
\]
for every $L\ge1$, but $\lambda_1(T_{r_L}),\lambda_2(T_{r_L})\to1$.
Moreover $r_L\to1$ in probability, $\mathbb E r_L^s\to2$, and
$\Theta=1$. Thus the finite-cutoff guarantee in
Theorem~\ref{thm:critical-tail} cannot include $\tau=1$ for $K\ge1$,
even under the exact total moment cap and a prescribed positive floor.
The same family satisfies the tail cap at every $K\ge1$ simultaneously;
it depends on $\rho$ and $b$, but not on $K$.
\end{theorem}
\begin{proof}
Use the undiluted product tilt $R_L$ from the proof of
Theorem~\ref{thm:critical-exact}. It satisfies
\[
 R_L>0,\qquad \mathbb E R_L^s=1,\qquad
 \mathbb E R_L^t\longrightarrow0\quad(0<t<s).
\]
Write $m_L=\mathbb E R_L\le1$. The same proof supplies coefficients
$\beta_L\to1$ such that $w_L=(1+\beta_LR_L)/Z_L$, with
$Z_L=1+\beta_Lm_L\to1$, has both leading eigenvalues tending to one.
Here $\beta_L$ denotes the balancing coefficient previously called
$a_L$; the new coefficient below is different.
The product definition gives additional information at its own mean.
For $\delta_L=dL^{-1/3}$ and the single-site normalizer $n(\delta_L)$,
$m_L=n(\delta_L)^{-L}>0$, so exactly
\[
 \frac{R_L}{m_L}=\prod_{i=1}^L(1+\delta_L x_i),\qquad
 \mathbb E(R_L/m_L)^t=
 \left(\frac{(1+\delta_L)^t+(1-\delta_L)^t}{2}\right)^L.
\]
Fix $0<t<1$. The logarithm of the bracket is an even smooth function
on any fixed compact interval $|\delta|\le d<1$. Taylor's theorem,
with bounded fourth derivative there, yields
\[
 \log\frac{(1+\delta)^t+(1-\delta)^t}{2}
 =\frac{t(t-1)}2\delta^2+O_{t,d}(\delta^4).
\]
After multiplication by $L$, the leading term tends to minus infinity
as $L^{1/3}$ and the remainder tends to zero. Markov's inequality gives
$\mathbb P(R_L>m_L)\le\mathbb E(R_L/m_L)^t\to0$.
This relative-mean estimate does not follow from $R_L\to0$ alone.
Set
\[
 p_L=\mathbb P(R_L>m_L)^{1/s},\quad a_L=1-p_L,\quad
 A_L=1-a_Lm_L,\quad q_L=A_L+a_LR_L.
\]
All three coefficients belong to $[0,1]$, $\mathbb E q_L=1$, and
$a_L,A_L\to1$, whence $q_L\to1$ in probability. The identity
$q_L-1=a_L(R_L-m_L)$ shows that $E_L=\{q_L>1\}$ equals
$\{R_L>m_L\}$ when $a_L>0$ and is empty when $a_L=0$.
In both cases the event inclusion holds. Minkowski gives
\[
 \|q_L\mathbf1_{E_L}\|_s
 \le A_L\mathbb P(E_L)^{1/s}+a_L\|R_L\mathbf1_{E_L}\|_s
 \le p_L+a_L=1.
\]
No division by $a_L$ occurs; its zero endpoint gives $q_L=1$.

Put $S(v)=QM_vQ$. Endpoint forward hypercontractivity maps $L^2$
to $L^{2s/(s-1)}$ with norm at most one. H\"older therefore gives
$\|QM_hQ\|\le\|h\|_s$ for every real $h$. In particular,
$\|QM_{R_L}Q\|\le1$, and
\[
 \|S(q_L)-S(w_L)\|
 \le |A_L-Z_L^{-1}|+|a_L-\beta_L/Z_L|\longrightarrow0.
\]
The square matrices $S(v)$ and $T_v$ are $X^*X$ and $XX^*$ for
$X=M_{\sqrt v}Q$, so are isospectral with multiplicities. The
min-max perturbation bound thus transfers both actual leading
eigenvalue limits to $T_{q_L}$.

We next prove $M_L:=\mathbb E q_L^s\to2$. The lower bound is
$M_L\ge A_L^s+a_L^s\mathbb E R_L^s\to2$. For a fixed $T>0$,
$q_L^s\mathbf1_{R_L\le T}$ converges in probability to one and is
bounded by $(1+T)^s$, so its expectation tends to one. On $R_L>T$,
$q_L^s\le(1+1/T)^sR_L^s$. Consequently
$\limsup_L M_L\le1+(1+1/T)^s$. Letting $T\to\infty$ proves the
claim. This auxiliary truncation is used only for the moment limit.

Set $\varepsilon_L=0$ when $M_L\le2$, and otherwise set
$\varepsilon_L=(M_L-2)/(M_L-1)$. Then $0\le\varepsilon_L<1$ and
$\varepsilon_L\to0$. The mean-one density
$v_L=(1-\varepsilon_L)q_L+\varepsilon_L$ satisfies, by convexity,
\[
 \mathbb E v_L^s\le(1-\varepsilon_L)M_L+\varepsilon_L\le2.
\]
Since $\varepsilon_L<1$, the events $v_L>1$ and $q_L>1$ coincide.
On this event $v_L\le q_L$, so pointwise
$v_L^s\mathbf1_{v_L>1}\le q_L^s\mathbf1_{q_L>1}$. The strict
tail cap at one is preserved. Also
$S(v_L)-S(q_L)=\varepsilon_L(P_\rho-S(q_L))\to0$ in operator norm,
since $\|S(q_L)\|\to1$. Both eigenvalue limits persist, as does
convergence in probability to one. The two-sign estimate
\eqref{eq:two-sign-budget} gives
$2\lambda_2(T_{v_L})^s\le\mathbb E v_L^s\le2$, and hence the
critical moment tends to two.

Since $v_L\ge A_L\to1$, replace a finite initial segment by $r_L=1$
and use $r_L=v_L$ thereafter to impose $r_L\ge b$ in every dimension.
The prefix has empty strict tail above one, and all limits are unchanged.
For every $K\ge1$, the tail integrand is pointwise at most the one
at cutoff one. Thus this same family works simultaneously for all such $K$.
For each fixed $J>1$, the bounded truncation
$r_L^s\mathbf1_{r_L\le J}$ has expectation tending to one. Thus the
$J$-tail tends to one, giving $H(J)\ge1$. The
all-dimension tail bound at one also gives $H(J)\le1$ for every $J>1$.
Therefore $H(J)=1$ for every $J>1$ and $\Theta=1$.
\end{proof}

Each finite matrix in this construction has a positive Perron gap;
the gaps tend to zero with dimension. The two moment caps are
inequalities in each dimension, with total moment equality asserted
only in the limit. The mean-relative probability coefficient resolves
the cutoff-one case without a dimension-dependent cutoff.


\begin{corollary}[Cutoffs below one]\label{cor:cutoff-below-one}
For $0\le K<1$, a nonnegative mean-one density satisfying
$\mathbb E[r^s\mathbf1_{r>K}]\le1$ has
\[
 \gap(T_r)\ge1-\left(\frac{1+K^{s-1}}2\right)^{1/s}>0.
\]
\end{corollary}
\begin{proof}
Pointwise $r^s\mathbf1_{r\le K}\le K^{s-1}r$, including $K=0$.
Thus $\mathbb E r^s\le1+K^{s-1}<2$. The two-sign bound and
$\lambda_1\ge1$ give
$\lambda_2/\lambda_1\le\lambda_2\le(\mathbb E r^s/2)^{1/s}$.
\end{proof}
No optimality is claimed for this lower constant or for a general
tail-budget threshold below cutoff one. The displayed bound is uniform
in dimension for fixed $K<1$, not uniform as $K$ approaches one.

These are worst-case statements about a specialized fixed Boolean
congruence. Scalar tails under the base measure are not the operator-tail
functionals under a Perron invariant measure in~\cite{Wang}. The cutoff
bound refines the same sign-truncation proof; it is not an independent
method or a claim of historical priority. The qualitative boundary is
sharp, while the positive constant is not optimized.

\section{Optimized finite backgrounds below cutoff one}\label{sec:background}
Fix $0<\rho<1$ and, throughout this section, set
\[
 s=\frac{1+\rho}{1-\rho},\qquad k_s=\frac{s-1}{s},\qquad
 C_s=\frac{s^s}{(s-1)^{s-1}}.
\]
The cutoff-one result does not classify a fixed cutoff $K<1$ with
a tail budget greater than one. We now combine long parities with the
balanced critical tilt to give a sub-two collapse region. The sharp
optimization below concerns a scalar background law, not the optimal
spectral boundary among all densities.

\begin{proposition}[Scalar background tail optimum]\label{prop:background-optimum}
For $s>1$ and $0<K<1$, among nonnegative random variables $A$ with
$\mathbb E A=1$,
\[
 \inf\mathbb E[A^s\mathbf1_{A>K}]=m_s(K),\qquad
 m_s(K)=\begin{cases}
 1,&0<K\le k_s,\\
 C_s(1-K)K^{s-1},&k_s<K<1.
 \end{cases}
\]
For $K\le k_s$ the constant law $A=1$ attains the infimum. For $K>k_s$
an attaining law has values $K$ and $v_*=sK/(s-1)$, with high probability
$p_*=(s-1)(1-K)/K$. Replacing the strict event by $A\ge K$ leaves
the infimum unchanged, but that infimum is not attained when $K>k_s$.
\end{proposition}
\begin{proof}
For $K\le k_s$, convexity gives $y^s\ge1+s(y-1)$. If $y\le K$, the
right side is nonpositive. Hence $y^s\mathbf1_{y>K}\ge1+s(y-1)$
everywhere. Taking expectations gives the lower bound one, attained
by $A=1$. The affine minorant is integrable even if the left side has
infinite expectation.

For $K>k_s$, the derivative of $y^s/(y-K)$ on $y>K$ has the sign
of $(s-1)y-sK$. Its unique minimum is at $v_*$ and its value is
$c=C_sK^{s-1}$. Thus
\[
 y^s\mathbf1_{y>K}\ge c(y-K)\qquad(y\ge0).
\]
Integration gives $c(1-K)$. The stated probabilities lie in $(0,1)$,
have mean $K+p_*(v_*-K)=1$, and put mass only at the two equality
points of the minorant. At $K=k_s$ this law degenerates to $A=1$.

For the non-strict event the same lower bound holds. When $K>k_s$,
choose $u<K$ tending to $K$ with $u>k_s$, high value $su/(s-1)$,
and high probability $(s-1)(1-u)/u$. The high value exceeds $K$
eventually, and the cost tends to $C_s(1-K)K^{s-1}$. Thus the
infimum is unchanged. Equality in the affine minorant for the
non-strict integrand is possible only at $v_*$: below $K$ its
right side is strictly negative, and at $K$ its left side is $K^s>0$.
Concentration at $v_*>1$ cannot have mean one, proving nonattainment.
For $K\le k_s$, $A=1$ still attains one.
\end{proof}

This is an elementary instance of classical moment optimization by
pointwise affine bounds; see Bertsimas and Popescu~\cite{BertsimasPopescu}
for the general moment-duality setting. Neither that method nor the
formula's external priority is claimed as new here.

\begin{theorem}[Sub-two tail budgets with finite parity backgrounds]
\label{thm:background-collapse}
Fix $k_s<K<1$, $0<b<K$, and
\[
 1+C_s(1-K)K^{s-1}<\tau<2.
\]
There are strictly positive densities $r_L$ on every Boolean cube
of integer dimension $L\ge1$ satisfying
\[
 \mathbb E r_L=1,\qquad r_L\ge b,\qquad
 \mathbb E[r_L^s\mathbf1_{r_L>K}]\le\tau,
\]
such that both leading eigenvalues of $T_{r_L}$ tend to one and
$\gap(r_L)\to0$. The densities use global interactions. No marginal
matching, locality, equality-budget attainability or optimal universal
spectral threshold is asserted.
\end{theorem}
\begin{proof}
\emph{A fixed dyadic background.}
Write $g(u)=C_s(1-u)u^{s-1}$. On $(k_s,1)$,
\[
 g'(u)=C_su^{s-2}((s-1)-su)<0,
\]
and its endpoint values are one and zero. Strict budget slack permits
a fixed $\max(b,k_s)<u<K$ with $1+g(u)<\tau$. The optimal high
probability at lower value $u$ is $p_u=(s-1)(1-u)/u$.
For $p\in(0,1)$ define $v(p)=u+(1-u)/p$. The mean
$(1-p)u+pv(p)$ is exactly one and $pv(p)^s$ is continuous in $p$.
Choose a fixed dyadic $p=j/2^k$ sufficiently close to $p_u$ that
\[
 1+pv^s<\tau,\qquad v=v(p)>1>K,
 \qquad k\ge1,\quad 1\le j<2^k.
\]
On $k$ unbiased signs choose a function $F$ equal to $v$ at $j$
patterns and equal to $u$ elsewhere. Then $\mathbb EF=1$.
Replace each sign by the parity of a disjoint block of $M$ unbiased
bits, obtaining $a_M$ on $kM$ bits. Those parities are independent
unbiased signs, so $a_M$ has exactly the same two-point law.
All $u,v,p,k$ are fixed before the dimension grows.

\emph{An operator-norm estimate for every character.}
Let $Q_d=P_{\sqrt\rho,d}$ and let $\chi_U$ be a Walsh character.
For any Walsh basis vector $\chi_S$,
\[
 Q_dM_{\chi_U}Q_d\chi_S
 =\rho^{(|S|+|S\mathbin\triangle U|)/2}\chi_{S\mathbin\triangle U}
 =\rho^{|U|/2+|S\setminus U|}\chi_{S\mathbin\triangle U}.
\]
This is a weighted permutation of an orthonormal basis. Its norm is
exactly $\rho^{|U|/2}$, attained for $S\subset U$.
Writing the finite Walsh expansion
$F=1+\sum_{I\ne\varnothing}\widehat F(I)\chi_I$ and putting
$D_F=\sum_{I\ne\varnothing}|\widehat F(I)|$, block substitution gives
\begin{equation}\label{eq:318-background-decay}
 \|Q_{kM}M_{a_M}Q_{kM}-P_{\rho,kM}\|
 \le \sum_{I\ne\varnothing}|\widehat F(I)|\rho^{M|I|/2}
 \le D_F\rho^{M/2}\longrightarrow0.
\end{equation}
Large fixed amplitudes or coefficients do not obstruct this limit.

\emph{The inherited critical family and actual eigenvalues.}
Use the balanced critical family in Theorem~\ref{thm:critical-exact}:
with fixed small $d>0$, $\delta_n=dn^{-1/3}$ and
\[
 N(\delta)=\left(\frac{(1+\delta)^s+(1-\delta)^s}{2}\right)^{1/s},
 \qquad R_n(y)=\prod_{i=1}^n\frac{1+\delta_n y_i}{N(\delta_n)},
\]
put
\[
 \kappa(\delta)=\frac{1+\rho+
 \sqrt{(1-\rho)^2+4\rho\delta^2}}{2N(\delta)},\quad
 \beta_n=\kappa(\delta_n)^{-n},\quad
 Z_n=1+\beta_n\mathbb ER_n,\quad
 w_n=\frac{1+\beta_nR_n}{Z_n}.
\]
The established balanced Gram estimates give the actual limits
$\lambda_1(T_{w_n}),\lambda_2(T_{w_n})\to1$, not merely two
trial-vector bounds. Also $\beta_n,Z_n\to1$, $\mathbb ER_n^s=1$,
and $\mathbb ER_n^t\to0$ for every fixed $0<t<s$.
The latter follows from the one-site expansion
\[
 \log\mathbb E\left[\frac{1+\delta Y}{N(\delta)}\right]^t
 =\frac{t(t-s)}2\delta^2+O(\delta^4),
\]
since $n\delta_n^2\to\infty$ and $n\delta_n^4\to0$.

For sufficiently large $L$ take $M=\lfloor L/(2k)\rfloor$ and
$n=L-kM$. Both tend to infinity. On the two independent blocks set
\[
 r_L(x,y)=\frac{a_M(x)+\beta_nR_n(y)}{Z_n}.
\]
Its mean is one and its minimum is at least $u/Z_n$, eventually
greater than $b$. The matrices $S_r=Q_LM_rQ_L$ and $T_r$ have the
same full spectrum by the $AA^*$ and $A^*A$ identity with
$A=M_{\sqrt r}Q_L$. Linearity and tensorization show
\[
 \left\|S_{r_L}-P_{\rho,kM}\otimes S_{w_n}\right\|
 \le D_F\rho^{M/2}/Z_n\longrightarrow0.
\]
The tensor comparison has largest eigenvalue $\lambda_1(S_{w_n})$
and second eigenvalue between $\lambda_2(S_{w_n})$ and
$\lambda_1(S_{w_n})$: use the constant background eigenvector
and the first two tilt eigenvectors for the lower bound. Positivity
and the fact that the top background eigenvalue is one give the
upper bound. Both leading eigenvalues therefore tend to one.
The self-adjoint eigenvalue perturbation inequality transfers both
limits to $S_{r_L}$ and hence $T_{r_L}$. Each finite matrix remains
strictly positive, with a simple Perron root.

\emph{Noninteger moments and the moving strict cutoff.}
Let $\theta=\min(1,s/2)$. For $c>0,z\ge0$,
\begin{equation}\label{eq:318-mixed-powers}
 0\le(c+z)^s-c^s-z^s
 \le s2^{s-1}\big(c^{s-\theta}z^\theta+
                   c^\theta z^{s-\theta}\big).
\end{equation}
To prove this, divide by $c^s$ and put $y=z/c$. For $y\le1$,
the mean value theorem bounds $(1+y)^s-1-y^s$ by
$s2^{s-1}y\le s2^{s-1}y^\theta$. For $y\ge1$ interchange
$c$ and $z$. Nonnegativity follows from convexity. Both powers
$\theta,s-\theta$ lie strictly between zero and $s$.
For $z=\beta_nR_n$, the mixed expectations vanish and
$\mathbb Ez^s\to1$. Thus $\mathbb E(c+\beta_nR_n)^s\to c^s+1$.

In the high sector, $v/Z_n>K$ eventually, so the entire sector
contributes and its conditional tail moment tends to $v^s+1$.
In the low sector the event is exactly $R_n>c_n$, where
$c_n=(KZ_n-u)/\beta_n\to K-u>0$. Fix positive finite
$c_-,c_+$ containing these thresholds eventually. For $0<t<s$,
\[
 \mathbb P(R_n>c_n)\le c_-^{-t}\mathbb ER_n^t\to0,
 \qquad
 \mathbb E[R_n^s\mathbf1_{R_n\le c_n}]
 \le c_+^{s-t}\mathbb ER_n^t\to0.
\]
Consequently the strict event retains asymptotic $s$-moment one.
In \eqref{eq:318-mixed-powers} restricted to this event, the constant term vanishes with
its probability, and the mixed terms are bounded by their vanishing
unrestricted expectations. After normalization, the low-sector
tail moment tends to one. Averaging the fixed sector probabilities,
\[
 \mathbb E[r_L^s\mathbf1_{r_L>K}]
 \longrightarrow (1-p)+p(v^s+1)=1+pv^s<\tau.
\]
This calculation does not infer moment convergence from convergence
in probability. Indeed the full moment tends to
$1+(1-p)u^s+pv^s>2$, by strict Jensen for the nonconstant background.

Choose a common $L_0$ beyond which both blocks are positive, the
floor holds and the tail cost is at most $\tau$. For $1\le L<L_0$
set $r_L=1$. Its tail cost is one and its floor is valid since
$b<K<1$. This finite replacement preserves all limits and proves
the assertion in every dimension. No effective value of $L_0$ is claimed.
\end{proof}

\subsection*{Comparison of sufficient and exact spectral thresholds}
The half-parity background with low value approaching $K$ and high
value approaching $2-K$ has limiting scalar cost $(2-K)^s/2$.
It gives collapse whenever
\[
 2-2^{1/s}<K<1,\qquad 0<b<K,\qquad
 1+(2-K)^s/2<\tau<2.
\]
This is a specialization of the same proof with $p=1/2$ and
$u<K$ chosen with strict slack. Proposition~\ref{prop:background-optimum}
shows that the optimized cost is no larger. Equality occurs only
when $2-K=sK/(s-1)$, that is $K=2(s-1)/(2s-1)$; outside this
isolated value the cost improvement is strict. Strict convexity
of $2^t$ below its chord on $0<t<1$ gives
$2^{1/s}<1+1/s$. Hence the new interval $k_s<K<1$ strictly
extends the half-parity sub-two interval.

On the other hand, any mean-one density with tail budget at most
$\tau$ obeys $\mathbb Er^s\le K^{s-1}+\tau$, since
$r^s\mathbf1_{r\le K}\le K^{s-1}r\mathbf1_{r\le K}$.
The established total-moment estimate therefore gives
\[
 \gap(r)\ge1-\left(\frac{K^{s-1}+\tau}{2}\right)^{1/s}>0
 \quad\hbox{if }\tau<2-K^{s-1}.
\]
This sufficient estimate alone does not identify the universal spectral
threshold. The construction justifies adding an escaped unit to the scalar
optimum only for that particular family. Theorem~\ref{thm:attained-cutoff}
attains this budget with floor $b=K$, while the separate necessary estimate
in Theorem~\ref{thm:universal-tail} proves its universal optimality.
For $K<b<1$ the event contains every point
and reduces to the total-moment condition. At $b=1$ mean one forces
$r=1$ and gap $1-\rho$; at $b>1$ the class is empty. The earlier
source309A omitted $b<1$ in its additional comparison; its failed
gate and the corrected309B report remain preserved, not relabelled.


\section{Attaining the constructive cutoff budget}\label{sec:attained-cutoff}
The preceding fixed-background argument uses strict slack in both the
floor and the tail budget. We now remove both restrictions. The cutoff
event remains strictly $r>K$; mass exactly at $K$ is essential.

\begin{theorem}\label{thm:attained-cutoff}
Fix $0<\rho<1$, $s=(1+\rho)/(1-\rho)$, and
\[
 k_s=\frac{s-1}{s}<K<1,\qquad
 g(K)=\frac{s^s}{(s-1)^{s-1}}(1-K)K^{s-1}.
\]
For every integer $L\ge1$ there is a strictly positive density $r_L$
on the uniform Boolean cube such that
\[
 \mathbb Er_L=1,\qquad r_L\ge K,\qquad
 \mathbb E[r_L^s\mathbf1_{r_L>K}]\le1+g(K),
\]
and $\lambda_1(T_{r_L})\to1$, $\lambda_2(T_{r_L})\to1$.
Consequently $\gap(T_{r_L})\to0$. The same family meets every
prescribed floor $0<b\le K$ and every budget $\tau\ge1+g(K)$.
This is an attained constructive upper budget, not a claim that
$1+g(K)$ is the optimal threshold over all collapsing families.
\end{theorem}

\begin{proof}
All expectations below use normalized uniform probability. Write
$Q_d=P_{\sqrt\rho,d}$. For positive $r$, the full spectra of
$T_r=M_{\sqrt r}P_{\rho,d}M_{\sqrt r}$ and $S_r=Q_dM_rQ_d$
agree, by the $AA^*/A^*A$ identity with $A=M_{\sqrt r}Q_d$.

\emph{The inherited balanced tilt.}
Use the explicit family in the proof of
Theorem~\ref{thm:critical-exact}. Fix a sufficiently small $d>0$ and set
\[
 \delta_n=dn^{-1/3},\qquad
 N(\delta)=\left(\frac{(1+\delta)^s+(1-\delta)^s}{2}\right)^{1/s},
 \qquad
 R_n(y)=\prod_{i=1}^n\left(\frac{1+\delta_n y_i}{N(\delta_n)}\right).
\]
In particular, the normalization is applied to every product factor.
Define
\[
 H_n=Q_nM_{R_n}Q_n,\qquad
 \kappa(\delta)=\frac{1+\rho+
 \sqrt{(1-\rho)^2+4\rho\delta^2}}{2N(\delta)},\qquad
 b_n=\kappa(\delta_n)^{-n},\qquad
 F_n=P_{\rho,n}+b_nH_n.
\]
The earlier proof gives
\begin{equation}\label{eq:attained-inherited}
 b_n\to1,\quad \|H_n\|=\kappa(\delta_n)^n\to1,\quad
 \lambda_1(F_n)\to1,\quad\lambda_2(F_n)\to1,
 \quad\mathbb ER_n^s=1.
\end{equation}
These are conclusions for that explicit family, not additional hypotheses.
For completeness, its interior moments have the exact expression
\[
 a_n(t):=\mathbb ER_n^t
 =\left(\frac{(1+\delta_n)^t+(1-\delta_n)^t}
 {2N(\delta_n)^t}\right)^n\longrightarrow0\qquad(0<t<s).
\]
Indeed their logarithm is
$n[t(t-s)\delta_n^2/2+O(\delta_n^4)]$.
At the critical exponent, $\log\kappa(\delta)=O(\delta^4)$.
The block Gram comparison proved earlier has two diagonal top
eigenvalues one and an off-diagonal norm tending to zero; it controls
the actual two leading eigenvalues of $F_n$. Its second lower bound
is eventually positive, so it is not an added zero of the Gram matrix.

\emph{Growing finite background laws.}
The scalar optimizer from Proposition~\ref{prop:background-optimum} has
\[
 v_*=\frac{sK}{s-1}>1,\qquad
 p_*=\frac{(s-1)(1-K)}K\in(0,1),\qquad p_*v_*^s=g(K).
\]
For sufficiently large $L$ choose
\[
 k=\lfloor L^{1/3}\rfloor,\quad
 M=\left\lfloor\frac{L}{2k}\right\rfloor,\quad n=L-kM,
 \qquad p=p_L=\frac{\lfloor2^kp_*\rfloor}{2^k},
 \qquad v=v_L=K+\frac{1-K}{p}.
\]
Then $p\in[p_*/2,p_*]$, $p\to p_*$, $v\to v_*$ and
$v\le V:=2v_*$. Also $k,M,n\to\infty$, $n\ge L/2$ and
$M/k\to\infty$. Choose exactly $2^kp$ patterns in $\{-1,1\}^k$,
and let $f_L$ take value $v$ there and $K$ elsewhere.
Replace its $k$ signs by the parities of $k$ disjoint blocks of $M$
bits. Call the resulting background $A_L$, with high-sector indicator
$I_L$. Independent unbiased parities give
\[
 \mathbb EI_L=p,\qquad A_L=K+(v-K)I_L,\qquad\mathbb EA_L=1.
\]
For a nonempty Walsh support $U$,
\[
 Q_dM_{\chi_U}Q_d\chi_S
 =\rho^{|U|/2+|S\setminus U|}\chi_{S\mathbin\triangle U}.
\]
Thus this weighted permutation has norm exactly $\rho^{|U|/2}$.
Each Walsh coefficient of $f_L$ has absolute value at most
$\mathbb E f_L=1$, and there are at most $2^k-1$ nonconstant ones.
Every substituted support has size at least $M$, so
\begin{equation}\label{eq:attained-parity}
 \|Q_{kM}M_{A_L}Q_{kM}-P_{\rho,kM}\|
 \le(2^k-1)\rho^{M/2}\longrightarrow0.
\end{equation}
The limit follows from $k\log2+(M/2)\log\rho\to-\infty$.
In particular, the growing number of coefficients is explicitly
controlled. The scalar affine minorant gives
\begin{equation}\label{eq:attained-dyadic-excess}
 d_L:=p_Lv_L^s-g(K)\ge0,\qquad d_L\to0.
\end{equation}
No uniform finite cube is assumed to realize a non-dyadic $p_*$ exactly.

\emph{Clipping, exact mean, and the amplitude reserve.}
Put $\theta=\min(1,s/2)$, so both $\theta$ and $s-\theta$ lie in $(0,s)$.
For $n=n_L$ define
\[
 c_n=a_n(\theta)^{1/(2\theta)},\qquad
 U_n=(R_n-c_n)_+,\qquad m'_n=\mathbb EU_n.
\]
Then $c_n\to0$, $m'_n\le a_n(1)\to0$, and
\begin{equation}\label{eq:attained-clipping}
 \mathbb P(U_n>0)\le\sqrt{a_n(\theta)},\qquad
 \|Q_nM_{U_n-R_n}Q_n\|\le c_n.
\end{equation}
The first inequality is Markov's inequality at $c_n$; the second uses
$0\le R_n-U_n\le c_n$ and $\|Q_n\|=1$.
Equality $R_n=c_n$ is in the complement of the strict event.

For real $s>1$, $c>0$ and $z\ge0$, the mixed-power bound is
\begin{equation}\label{eq:attained-mixed}
 0\le(c+z)^s-c^s-z^s
 \le B_s\bigl(c^{s-\theta}z^\theta+c^\theta z^{s-\theta}\bigr),
 \qquad B_s=s2^{s-1}.
\end{equation}
To see the upper bound, divide by $c^s$. For $z/c\le1$ the
mean value theorem bounds the difference by $B_s(z/c)$, which is
at most $B_s(z/c)^\theta$. For $z/c\ge1$ exchange $c,z$.
The lower bound follows by integrating a positive derivative.
Thus this estimate does not require an integer exponent.
Set
\begin{align*}
 E_L={}&K^s\sqrt{a_n(\theta)}\\
 &+B_s\bigl[(K^{s-\theta}+V^{s-\theta})a_n(\theta)
 +(K^\theta+V^\theta)a_n(s-\theta)\bigr],\\
 e_L={}&d_L+E_L.
\end{align*}
Both are nonnegative and tend to zero. Eventually $e_L<1/2$ and
$m'_n<1-K$. On that range let
\begin{equation}\label{eq:attained-density}
 \beta_L=(1-e_L)^{1/s},\qquad
 \gamma_L=\frac{\beta_Lm'_n}{p_L},\qquad
 r_L(x,y)=A_L(x)-\gamma_L I_L(x)+\beta_LU_n(y).
\end{equation}
Here $0<\beta_L\le1$, $\beta_L\to1$ and $\gamma_L\to0$.
The mean is exactly $1-p_L\gamma_L+\beta_Lm'_n=1$.
The low sector has density $K+\beta_LU_n\ge K$; on the high sector,
\[
 v_L-\gamma_L=K+\frac{1-K-\beta_Lm'_n}{p_L}>K.
\]
Thus mean compensation has not lowered the floor. It acts only
on the high sector, and no global normalization is applied.

On the low sector the strict event is precisely $U_n>0$. Apply
\eqref{eq:attained-mixed} with $c=K$ and $z=\beta_LU_n$, restricted
to that event. Its constant contribution is at most
$K^s\mathbb P(U_n>0)$. The high sector is entirely above $K$;
discarding the nonnegative subtraction $\gamma_L$ gives an upper
bound from \eqref{eq:attained-mixed} with $c=v_L$.
Averaging these bounds with probabilities $1-p_L,p_L$, the
coefficient of $\beta_L^s\mathbb EU_n^s$ is exactly one.
Use $U_n\le R_n$, $v_L\le V$, $\beta_L\le1$,
$\mathbb ER_n^s=1$, and \eqref{eq:attained-clipping}. The two mixed
terms are bounded by those defining $E_L$. Hence
\begin{equation}\label{eq:attained-exact-budget}
 \mathbb E[r_L^s\mathbf1_{r_L>K}]
 \le p_Lv_L^s+\beta_L^s+E_L
 =g(K)+d_L+1-e_L+E_L=1+g(K).
\end{equation}
The reserve pays for the dyadic excess and every retained tail error.
No convergence of $s$-moments was inferred from convergence in probability.

\emph{The actual leading eigenvalues.}
Write $P_b=P_{\rho,kM}$, $P_t=P_{\rho,n}$ and similarly $Q_b,Q_t$.
By linearity and tensorization,
\[
 S_{r_L}=(Q_bM_{A_L}Q_b-\gamma_LQ_bM_{I_L}Q_b)\otimes P_t
 +\beta_LP_b\otimes(Q_tM_{U_n}Q_t).
\]
Compare with $B_L=P_b\otimes F_n$. Since $0\le I_L\le1$,
\eqref{eq:attained-parity} and \eqref{eq:attained-clipping} imply
\begin{align*}
 \|S_{r_L}-B_L\|
 &\le(2^k-1)\rho^{M/2}+\gamma_L
       +\beta_Lc_n+|\beta_L-b_n|\|H_n\|\longrightarrow0.
\end{align*}
All four errors vanish by the established limits.
Both tensor factors of $B_L$ are positive definite and
$\lambda_1(P_b)=1$. Therefore
\[
 \lambda_1(B_L)=\lambda_1(F_n),\qquad
 \lambda_2(F_n)\le\lambda_2(B_L)\le\lambda_1(F_n).
\]
The lower bound uses two actual product eigenvectors with constant
background factor. In view of \eqref{eq:attained-inherited}, both
leading eigenvalues of $B_L$ tend to one. Weyl's self-adjoint
perturbation bound transfers both limits to $S_{r_L}$, and the
full-spectrum identity transfers them to $T_{r_L}$. This controls
the top eigenvalue as well as the second; two trial vectors alone
would not suffice.

Choose one $L_0$ so that every eventual condition above holds for
all $L\ge L_0$. For $1\le L<L_0$ set $r_L=1$. Its floor is at
least $K$, its mean is one, and its strict $K$-tail cost is
$1\le1+g(K)$. This supplies every dimension without changing the
limits. No effective estimate of $L_0$ is asserted.
\end{proof}

For the non-strict event $r_L\ge K$, every point contributes and the
tail becomes the total moment. The theorem therefore cannot be
relabeled with a non-strict event. The scalar optimization, clipping
and perturbation tools are classical; this construction and its
source audit do not establish external priority. The next two sections
improve the sufficient estimate and prove that the constructed upper
budget is the exact universal strict-tail threshold.
No marginal matching, locality, heat-bath gap or physical conclusion
is added to the separate families proved earlier.


\section{The exact scalar envelope at a strict cutoff}\label{sec:scalar-envelope}
This section sharpens the elementary sufficient estimate above. It is
useful for explicit bounds, but scalar optimality alone does not settle
the spectral question. Section~\ref{sec:universal-tail} supplies that step.
Throughout these two sections fix $0<\rho<1$, $s=(1+\rho)/(1-\rho)$,
$k_s=(s-1)/s<K<1$, and put
\[
 a=1-K,\quad p_*=(s-1)a/K,\quad C_s=\frac{s^s}{(s-1)^{s-1}},
 \quad g=C_saK^{s-1},\quad h(p)=\frac{(a+Kp)^s}{p^{s-1}}.
\]
Then $0<p_*<1$ and $0<g<1$. For $\tau\ge g$ let $p_\tau$ be the
unique solution of $h(p_\tau)=\tau$ in $(0,p_*]$, and set
$U(\tau)=\tau+K^s(1-p_\tau)$.

\begin{theorem}\label{thm:scalar-envelope}
Among nonnegative random variables on arbitrary probability spaces with
$\mathbb Er=1$ and $\mathbb E[r^s\mathbf1_{r>K}]\le\tau$, the class is
empty for $\tau<g$. For $\tau\ge g$ the maximum of $\mathbb Er^s$ is
$U(\tau)$, attained by the law with value $K$ of probability $1-p_\tau$
and value $K+a/p_\tau$ of probability $p_\tau$.
At $\tau=g$ this is the unique admissible law; a uniform $L$-cube realizes
it precisely when $2^Lp_*$ is integral. For $\tau>g$, the supremum over
all finite uniform Boolean cubes is $U(\tau)$, approached by positive
mean-one densities with floor $K$ and strict tail at most $\tau$.
\end{theorem}
\begin{proof}
Logarithmic differentiation gives
\[
 \frac{h'(p)}{h(p)}=\frac{Kp-(s-1)a}{p(a+Kp)}.
\]
Thus $h$ decreases from infinity to $h(p_*)=g$ on $(0,p_*]$, then
increases to $h(1)=1$. For an admissible $r$, write
$p=\mathbb P(r>K)>0$ and $m=\mathbb E[r\mathbf1_{r>K}]$.
The complementary mean is at most $K(1-p)$, so $m\ge a+Kp$.
Conditional Jensen yields
\[
 \tau\ge\mathbb E[r^s\mathbf1_{r>K}]
 \ge m^s/p^{s-1}\ge h(p)\ge g.
\]
It also forces $p\ge p_\tau$. Consequently
\[
 \mathbb Er^s\le\tau+K^s(1-p)\le U(\tau).
\]
The stated two-point law has mean one, strict tail $h(p_\tau)=\tau$
and total moment $U(\tau)$, proving attainment. If $\tau=g$, every
inequality above is equality: $p=p_*$, the lower value is almost surely
$K$, and strict convexity forces a constant upper value. This proves
uniqueness and the finite-cube integrality condition.
For $\tau>g$, choose
$p_L=(\lfloor2^Lp_\tau\rfloor+1)/2^L$. Eventually
$p_\tau<p_L<p_*$ and $p_L\to p_\tau$. Assign value $K+a/p_L$ to
exactly $2^Lp_L$ vertices and $K$ to the others. Its mean is one and its
strict tail $h(p_L)<\tau$, whereas its total moment tends to $U(\tau)$.
Exact attainment on a specified cube requires the high probability
$p_\tau$ to be realizable there; no irrational probability is assigned
to a finite uniform event.
\end{proof}

The earlier sign compression, $2\lambda_2^s\le\mathbb Er^s$ and
$\lambda_1\ge1$, now gives
\[
 \gap(T_r)\ge1-[U(\tau)/2]^{1/s}>0\qquad(U(\tau)<2).
\]
There is a unique $q_0\in(0,p_*)$ with
$h(q_0)+K^s(1-q_0)=2$, since this expression decreases strictly from
infinity to a value below two. Its tail budget
$\tau_c=2-K^s(1-q_0)$ satisfies
$2-K^{s-1}<2-K^s<\tau_c<2$; $U(\tau)<2$ exactly when
$g\le\tau<\tau_c$. For $s=2$ direct expansion gives
\[
 q_0=\frac{(1-K)^2}{2-2K+K^2},\qquad
 \tau_c=2-\frac{K^2}{2-2K+K^2}.
\]
In particular $K=3/4$ gives $\tau_c=25/17$, improving the elementary
$5/4$. This is only a sufficient spectral bound: a scalar extremizer
with total moment two need not saturate the two eigenvalue inequalities.
The moment optimization is an application of classical conditional
Jensen and extremal-law reasoning; see~\cite{BertsimasPopescu} for the
general method, without an assertion of priority for this formula.

\section{The universal strict-tail threshold}\label{sec:universal-tail}
The missing link is a necessary escaping-moment estimate for an arbitrary
collapsing family. Reverse hypercontractivity forces one normalized sign
part of a second eigenvector to have a small low quasi-norm. Truncation
then forces its energy into the high-density region. This proves the
additional unit of tail cost, without assuming the constructed density
floor or an a priori decomposition into background and tilt.

\begin{theorem}\label{thm:universal-tail}
Fix the parameters above. For every $g\le\tau<1+g$ there is an explicit
constant $\gamma_0>0$, depending only on $\rho,K,\tau$, such that every
nonnegative mean-one density on every finite uniform Boolean cube with
$\mathbb E[r^s\mathbf1_{r>K}]\le\tau$ satisfies
$\gap(T_r)>\gamma_0$. No positive lower bound on $r$ is required.
For $\tau<g$ the class is empty. At $\tau=1+g$, and hence at every larger
budget, Theorem~\ref{thm:attained-cutoff} gives a positive family with
floor $K$ and gap tending to zero. Thus $1+g$ is the exact universal
strict-tail collapse threshold for $k_s<K<1$.
\end{theorem}
\begin{proof}
Write $B=\tau+K^s$, $p=1-\rho$, $q=1+\rho$ and
$\theta=1/s^2$. Classical forward and reverse Boolean
hypercontractivity~\cite[Theorem~3.1 and Corollary~3.3]{MORSS} give
\begin{align}
 \langle F,P_\rho F\rangle
   &=\|P_{\sqrt\rho}F\|_2^2\le\|F\|_q^2,\label{eq:317-forward}\\
 \langle F,P_\rho G\rangle&\ge\|F\|_p\|G\|_p
                 \quad(F,G\ge0).\label{eq:317-reverse}
\end{align}
The reverse inequality has parameter condition
$\rho\le\sqrt{(1-p)(1-p)}$, with equality here, and permits zeros.
The quantities at exponent $p<1$ are quasi-norms; no triangle inequality
for them will be used. H\"older with $q/(2-q)=s$ gives
\[
 \|\sqrt z\,h\|_q^2\le(\mathbb Ez^s)^{1/s}\|h\|_2^2
 \quad(z\ge0).
\]
Since $\mathbb Er^s\le B$, this implies $\lambda_1\le B^{1/s}$.
The unit vector $\sqrt r$ has Rayleigh quotient
$\langle r,P_\rho r\rangle\ge(\mathbb Er)^2=1$, so
$\lambda_1\ge1$.

Put $\gamma=1-\lambda_2/\lambda_1$. If $\lambda_2=0$, the gap is one
and the assertion is immediate. Otherwise work on the active set
$\{r>0\}$. The restricted matrix has strictly positive entries and a
strictly positive Perron vector. A normalized real second eigenvector
vanishes off this set, is orthogonal to the Perron vector, and therefore
has both signs. Write it as $v=f_+-f_-$, with disjoint nonnegative parts,
$a=\|f_+\|_2>0$, $b=\|f_-\|_2>0$, $a^2+b^2=1$, and
$c=\langle f_+,T_rf_-\rangle>0$. The eigenvector equation implies
\[
 \langle f_+,T_rf_+\rangle=\lambda_2a^2+c,\qquad
 \langle f_-,T_rf_-\rangle=\lambda_2b^2+c.
\]
The compression to the orthonormal pair $u=f_+/a,w=f_-/b$ is
\[
 \lambda_2I+c\begin{pmatrix}a^{-2}&(ab)^{-1}\\(ab)^{-1}&b^{-2}\end{pmatrix}.
\]
Its upper eigenvalue is $\lambda_2+c/(a^2b^2)\le\lambda_1$.
Thus the \emph{normalized} cross term obeys
\[
 d=\langle u,T_rw\rangle=c/(ab)
       \le(\lambda_1-\lambda_2)ab\le B^{1/s}\gamma/2.
\]
Both sign Rayleigh quotients are at least $\lambda_2$.
Apply~\eqref{eq:317-reverse} to $\sqrt r\,u$ and $\sqrt r\,w$.
Choose the sign vector $h$ with the smaller $p$ quasi-norm after
multiplication by $\sqrt r$. For $F=\sqrt r\,h$ we obtain
\begin{equation}\label{eq:317-sign}
 \|F\|_p\le(B^{1/s}\gamma/2)^{1/2},\qquad
 \langle F,P_\rho F\rangle\ge\lambda_2\ge1-\gamma.
\end{equation}
The argument tolerates arbitrarily unbalanced sign norms $a,b$.

Fix $M>K$, let $F_0=F\mathbf1_{r\le M}$,
$F_1=F\mathbf1_{r>M}$, and $H_M=\mathbb E[r^s\mathbf1_{r>M}]$.
One checks $1/q=\theta/p+(1-\theta)/2$. Ordinary H\"older applied to
$Z^q=(Z^p)^{\theta q/p}(Z^2)^{(1-\theta)q/2}$ proves
\[
 \|Z\|_q\le\|Z\|_p^\theta\|Z\|_2^{1-\theta}.
\]
The two powers are positive and sum to one, so this use of H\"older has
exponents greater than one even though $p<1$.
Since $\|F_0\|_2\le\sqrt M$ and $\|F_0\|_p\le\|F\|_p$,
\eqref{eq:317-forward} and~\eqref{eq:317-sign} yield
\[
 E_0:=\langle F_0,P_\rho F_0\rangle
       \le M^{1-\theta}(B^{1/s}\gamma/2)^\theta.
\]
The same forward estimate and H\"older yield
$E_1:=\langle F_1,P_\rho F_1\rangle\le H_M^{1/s}$.
The Hilbert-space triangle after applying $P_{\sqrt\rho}$ gives
$\sqrt{\langle F,P_\rho F\rangle}\le\sqrt{E_0}+\sqrt{E_1}$.
Consequently
\begin{equation}\label{eq:317-escape}
 H_M\ge\left[\sqrt{1-\gamma}
 -M^{(1-\theta)/2}(B^{1/s}\gamma/2)^{\theta/2}\right]_+^{2s}.
\end{equation}
For a collapsing family with fixed $B$, this implies
$\liminf H_M\ge1$ for each fixed $M$.

Put $A=C_sK^{s-1}$, so $g=A(1-K)$. The pointwise scalar inequality
\[
 z^s\mathbf1_{z>K}\ge A(z-K)\qquad(z\ge0)
\]
holds because its right side is nonpositive below $K$, whereas
$z^s/(z-K)$ has its global minimum $A$ at $z=sK/(s-1)$ above $K$.
Apply it to $z=r\mathbf1_{r\le M}$ on the original probability space.
Since $\mathbb E[r\mathbf1_{r>M}]\le B/M^{s-1}$, the middle strict
tail costs at least $g-AB/M^{s-1}$. It is disjoint from the high tail,
so
\begin{equation}\label{eq:317-budget}
 \tau\ge H_M+g-AB/M^{s-1}.
\end{equation}
For any collapsing family, first fix $M$, use~\eqref{eq:317-escape},
and then send $M\to\infty$ in~\eqref{eq:317-budget}. This gives
$\tau\ge1+g$. No density floor or exchange of moment limits is involved.

For an explicit positive bound, let $\varepsilon=1+g-\tau\in(0,1]$ and
\begin{align*}
 M&=\max\{K+1,(4AB/\varepsilon)^{1/(s-1)}\},\\
 D&=1-(1-3\varepsilon/4)^{1/(2s)},\\
 \gamma_0&=\min\{D/4,\,
 2B^{-1/s}(D/4)^{2/\theta}M^{-(1-\theta)/\theta}\}.
\end{align*}
These constants are finite and positive, with $\gamma_0<1$.
If $\gamma\le\gamma_0$, the subtracted core term in
\eqref{eq:317-escape} is at most $D/4$ and
$\sqrt{1-\gamma}\ge1-\gamma\ge1-D/4$. Hence
\[
 H_M\ge(1-D/2)^{2s}>(1-D)^{2s}=1-3\varepsilon/4.
\]
But~\eqref{eq:317-budget} gives
$H_M\le\tau-g+\varepsilon/4=1-3\varepsilon/4$, a contradiction.
Thus $\gamma>\gamma_0$. The empty regime $\tau<g$ was proved in
Theorem~\ref{thm:scalar-envelope}; equality $\tau=g$ is included whenever
the finite-cube class exists. Finally Theorem~\ref{thm:attained-cutoff}
uses exactly this $\rho,s,K$ and strict event, and supplies collapse at
$1+g$ with floor $K$. This completes the sharpness assertion.
\end{proof}

For $s=2$, $\rho=1/3$ and $K=3/4$, the exact threshold is
$1+4(1/4)(3/4)=7/4$. The intermediate scalar bound $25/17$ remains
valid and can be numerically much simpler than the conservative explicit
constant above. Neither scalar optimality nor the construction alone
proved the universal threshold: the necessary escape
estimate~\eqref{eq:317-escape} closes that gap. External priority of this
specific synthesis is unestablished. Its inputs are classical; related
transmission of log-Sobolev and spectral-gap inequalities for perturbed
diffusion measures is studied, for example, in~\cite{Cattiaux}.
Those generator/measure results are not silently identified with the
finite transfer-matrix statement here. This theorem does not add marginal
matching, finite-range interactions or a physical interpretation.

\section{Comparison with hyperbounded and perturbation results}
The value $s_c$ is the classical Boolean hypercontractive boundary,
not a new integrability scale. Indeed, for $a_s=2s/(s+1)$, the
norm-one estimate $\|P_{\sqrt\rho}f\|_2\le\|f\|_{a_s}$ holds
when $s\ge s_c$ by the standard inequality recalled in~\cite{MOS}.
Its local sharpness is seen from $f=1+ux$ on one bit:
$\|P_{\sqrt\rho}f\|_2^2=1+\rho u^2$, whereas
$\|f\|_{a_s}^2=1+(a_s-1)u^2+o(u^2)$.
Tensorization amplifies a violated one-bit norm bound. Such top-norm
growth alone does not imply a closing normalized spectral gap;
Theorem~\ref{thm:tilts} adds two antipodal channels and their Gram
comparison to obtain a near-degenerate second eigenvalue.

The inequalities and matrix arguments used above are classical. Related
literature already supplies both qualitative gap criteria and, under
different hypotheses, dimension-independent quantitative perturbation
bounds. Miclo's Theorem~1 in~\cite{Miclo} gives a gap for ergodic
self-adjoint hyperbounded Markov operators. Section~6 and
Proposition~11 also discuss a norm-budget transition at two for the
fourth power of an $L^2$-to-$L^4$ Markov norm; this is an important
precedent for budget thresholds, not the raw density norm used here.
Hyperboundedness alone does not give a general quantitative gap. Wang's Theorem~1.1 in~\cite{Wang} gives a tail-norm criterion
for the gap of an invariant Markov operator's symmetrization.
These fixed-space conclusions must not be confused with a bound uniform
over a family of increasing dimensions.

Here $T_r$ is generally not Markov in $u_L$. Its Perron Doob transform
is Markov in its own stationary law, which generally differs from $r u_L$.
The ground-state multiplication is an $L^2$ isometry between these spaces,
but does not automatically identify their higher norms uniformly in $L$.
Thus an implication from a qualitative hyperboundedness theorem would
still need a quantitative change-of-measure argument.

Gross's Theorem~2.2 in~\cite{Gross} does give dimension-independent gaps
for additive Schr\"odinger perturbations on Riemannian probability
manifolds with a base logarithmic Sobolev inequality and suitable moments
of both $e^V$ and $e^{-V}$. We have not identified that additive operator
with the logarithm of this discrete transfer congruence. Our moment
hypothesis in Theorem~\ref{thm:moment} supplies no inverse-density moment
and permits zeros. In contrast, the counterexamples in
Theorems~\ref{thm:tilts} and~\ref{thm:critical-large} have $r\ge1/B$ and therefore bounded inverse
moments. For these families, missing inverse control is not a distinction
from Gross; the operator/form and exponent correspondences are the
unestablished steps.
Gl\"uck~\cite{Gluck}, especially Theorem~5.1, gives hyperboundedness criteria and uniform essential
spectral-radius estimates; those estimates alone do not control the
second eigenvalue of a finite matrix. These differences do not prove
originality or exclude a stronger implication. Aida's directly accessible author preprint~\cite{Aida}
uses an additive Schr\"odinger form with a derivation and chain-rule
hypotheses, together with ground-state and weak-Poincar\'e inputs.
Its setup and Theorem~10 were read in a targeted comparison; no
identification of those hypotheses with this multiplicative Boolean
transfer congruence has been established. No external priority, improvement over this literature, or exhaustive
novelty clearance is asserted.

\section{Relation to predecessors and reproducibility}
The congruence analysis in~\cite{ErikssonCongruence} permits positive
weights concentrating on an antipodal pair, with a transfer ratio tending
to one as volume grows. Such a pair fails to preserve pair marginals.
The parity analysis in~\cite{ErikssonParity} preserves every proper
marginal but concerns a heat-bath generator; Proposition~\ref{prop:parity}
shows precisely why it does not supply the transfer conclusion.
The present synthesis obtains linear-order marginal matching for both
the weight and the Perron stationary law, together with exponentially
many slow transfer modes.

The code ingredients and the matrix methods are classical. In particular,
Lemma~\ref{lem:code} is an elementary existence argument, not a new code
construction and not an efficient encoding claim. The citation to~\cite{AGM} supplies general limited-independence background, not a direct historical attribution for the dual-distance projection equivalence proved here. Work on simultaneously
good primal and dual codes, including~\cite{BrehmResch}, is stronger in
algorithmic aspects. Our targeted primary-source check did not identify
the exact joint conclusion above; this is not proof of historical priority.
External originality and the importance of this specialization remain
unestablished.

All proofs are finite-dimensional and analytical. No numerical experiment
or formal verification is needed or claimed. A source package retains the
preregistered proof checks, complete analytical audit, source hashes and
this manuscript. The permanent source tree is
\href{https://github.com/lluiseriksson/THE-ERIKSSON-PROGRAMME/tree/\SourceCommit/papers/marginal-transfer}{\texttt{papers/marginal-transfer} at the source commit below}.
The earlier spike proof is retained there in
\path{evidence296/finite-moment296A.txt}; the balanced critical proof and
its complete source audit are retained in \texttt{evidence298}.
The exact-budget proof, registered checks and full analytical source
audit are retained in \texttt{evidence299}. The full analytical reports,
contexts, assignments and source bindings for the grouped tail results are
in \texttt{evidence305} and \texttt{evidence307}. Their public-subset manifests explicitly identify
private reference PDFs and an unrelated local directory index omitted from
the original audit manifests; no redistribution of these files is implied.
The exact v4 PDF assessment is preserved there solely as historical evidence,
with all seven objections and a separate author response. Historical partial results
and exact-PDF reviews keep their original artifact identities; this
revision neither changes those reports nor inherits their grades.
The new background construction is grouped with these results because it
uses the same transfer operator, critical exponent and balanced tilt.
The complete source309A failure, corrected309B audit and source310 audit
are retained in \texttt{evidence311}, with their original contexts and
manifests. Three full third-party source captures from audit310 remain
private and are explicitly excluded from the public subset; their
original hashes are retained without implying redistribution rights.
The original v5 PDF and its complete first exact-PDF assessment, including
all six reservations, are historical evidence only. The expanded scalar
attainment discussion follows the recorded source310 audit clarification.
No v5 score is transferred to this manuscript.
The attained-cutoff construction is retained as source313, with its
complete analytical report, assignment, context and byte manifests in
\texttt{evidence314}. It removes two restrictions of the preceding
construction by clipping the tilt, compensating its mean only on the
high sector, and reserving amplitude for finite tail errors.
The source311 document and production-runner audits are also retained;
that intermediate manuscript was not compiled or given a PDF score.
A redundant complete frozen311 source archive remains private on D;
the public subset declares its exact hash and original Git source, while
retaining the full reports and original manifest. No report is shortened.
The sharp scalar envelope (source315) and reverse-hypercontractive necessity
argument (source317) are grouped here with the attained construction.
Their complete source reports, contexts and fixed judges are retained in
\texttt{evidence318}. The earlier exact v6 PDF, its full first-round
assessment and all six reservations are historical evidence, not a rating
of this revision. Its two explicitly forced equation numbers are replaced
by ordinary sequential numbering; the corresponding mathematical displays
are unchanged. The previous claim that the universal strict-tail interval
was unresolved is superseded by Theorem~\ref{thm:universal-tail}.
Document reproduction is a separate check of exact bytes
and layout, not an independent proof of the theorems.

% Consolidation322 draft. Sources319A and321C; not an assessed PDF.
\section{Growing marginal order at the critical moment budget}
\label{sec:growing-marginals}

The preceding marginal and moment constructions need not describe the same
weights. We now impose both requirements on one weighting law. Throughout
this section, $u_L$ is uniform probability on $\{-1,1\}^L$,
$0<\rho<1$ is fixed, and
\[
 s=\frac{1+\rho}{1-\rho},\qquad B_*=2^{1/s},\qquad
 Q=P_{\sqrt\rho},\qquad
 S_r=QM_rQ,\quad T_r=M_{\sqrt r}P_\rho M_{\sqrt r}.
\]
Here $P_t\chi_A=t^{|A|}\chi_A$. For $r\geq0$ with $\mathbb E r=1$,
$S_r$ and $T_r$ have the same eigenvalues, including zero multiplicities:
they are $AA^*$ and $A^*A$ for the square matrix $A=QM_{\sqrt r}$.
Also $\langle1,S_r1\rangle=1$, so $\lambda_1(T_r)\geq1$.
Exact uniform marginals of the weighting law $r u_L$ through order $k$
mean $\mathbb E[r\chi_A]=0$ for $1\leq|A|\leq k$.
This classical Fourier characterization and low-coefficient correction are prior methods\cite{alon-kwise,rubinfeld-xie}. The statement here concerns the weighting law;
no Perron stationary marginal assertion is made in this section.

Two different regimes result. At every fixed power $L^\alpha$ with
$\alpha<1/2$, exact matching is compatible with collapse at $B_*$.
At every fixed positive linear fraction of $L$, a bounded critical moment
instead forces convergence to the unweighted noise spectrum. The proof of
the latter statement uses the homogeneous moment inequality of
Kirshner and Samorodnitsky~\cite{kirshner-samorodnitsky}, not a new
hypercontractive inequality. The intervening sublinear regime is open here.

\begin{theorem}[Collapse with growing exact marginals]
\label{thm:growing-collapse}
Fix $0<\alpha<1/2$ and $0<b<1$. There are mean-one densities $r_L$,
for every integer $L\geq1$, such that
\[
 r_L\geq b,\qquad \|r_L\|_s\leq B_*,\qquad
 \mathbb E[r_L\chi_A]=0\quad(1\leq |A|\leq\lfloor L^\alpha\rfloor),
\]
while $\lambda_1(T_{r_L})\to1$, $\lambda_2(T_{r_L})\to1$ and
$\gap(T_{r_L})\to0$. Thus growing matching of these orders does not
raise the worst-case critical norm boundary $B_*$.
\end{theorem}

\begin{proof}
Choose $a=3/8-\alpha/4$, so $1/4<a<(1-\alpha)/2$, and fix a
sufficiently small $d\in(0,1)$. Set $\delta=dL^{-a}$ and
\[
 n(\delta)=\left(\frac{(1+\delta)^s+(1-\delta)^s}{2}\right)^{1/s},
 \quad R_L(x)=\prod_{i=1}^L\frac{1+\delta x_i}{n(\delta)},
 \quad m_L=\mathbb E R_L=n(\delta)^{-L}.
\]
The one-coordinate matrix and its largest eigenvalue are
\[
 H(\delta)=\frac1{n(\delta)}
 \begin{pmatrix}1&\delta\sqrt\rho\\\delta\sqrt\rho&\rho\end{pmatrix},
 \qquad
 \kappa(\delta)=\frac{1+\rho+
 \sqrt{(1-\rho)^2+4\rho\delta^2}}{2n(\delta)}.
\]
Writing $D=\diag(1,\rho)$ and
$\tau(\delta)=\|D^{1/2}H(\delta)^{1/2}\|$, direct diagonalization gives
\[
 \tau(\delta)^2=\frac{1+\rho^2+
 \sqrt{(1-\rho^2)^2+4\rho^2\delta^2}}{2n(\delta)}.
\]
At the fixed critical exponent the Taylor expansions are
\[
 \log n(\delta)=\frac{s-1}{2}\delta^2+O(\delta^4),\quad
 \log\kappa(\delta)=O(\delta^4),\quad
 \log\tau(\delta)=-\frac{\rho\delta^2}{2(1-\rho^2)}+O(\delta^4).
\]
Here $L\delta^4=d^4L^{\alpha-1/2}\to0$, whereas
$L\delta^2=d^2L^{1/4+\alpha/2}\to\infty$.
Consequently $c_L=\kappa(\delta)^{-L}\to1$, $m_L\to0$, and
$Z_L=1+c_Lm_L\to1$.

For $v_L=(1+c_LR_L)/Z_L$, the unnormalized operator is
$F_L=P_\rho+c_LH(\delta)^{\otimes L}$. Both summands have top
eigenvalue one. Form the Gram matrix of their positive square roots.
Its diagonal blocks are these two summands and its cross-block norm is
$\eta_L=\sqrt{c_L}\tau(\delta)^L\to0$. The block diagonal matrix has
its first two eigenvalues exactly one. Operator-norm perturbation puts
the first two Gram eigenvalues in $[1-\eta_L,1+\eta_L]$.
For large $L$ these are positive; the nonzero Gram spectrum is the spectrum
of $F_L$. Division by $Z_L$ therefore proves that the actual first two
eigenvalues of $S_{v_L}$ tend to one.

We also have $\mathbb E R_L^s=1$. For each fixed $0<t<s$,
\[
 \log\mathbb E R_L^t=-\frac{t(s-t)}2L\delta^2+O(L\delta^4)
 \longrightarrow-\infty.
\]
With $\theta=\min(1,s/2)$, the mean-value theorem for $y\leq1$ and
reciprocal substitution for $y\geq1$ give
\[
 0\leq(1+y)^s-1-y^s\leq
 s2^{s-1}(y^\theta+y^{s-\theta})\quad(y\geq0).
\]
Both exponents lie strictly between zero and $s$. Applying this inequality
to $c_LR_L$ yields $\mathbb E(1+c_LR_L)^s\to2$, hence
$\|v_L\|_s\to B_*$.

Put $k_L=\lfloor L^\alpha\rfloor$ and erase the nonconstant low levels:
\[
 p_L=m_L\sum_{1\leq|A|\leq k_L}\delta^{|A|}\chi_A,\qquad
 w_L=\frac{1+c_L(R_L-p_L)}{Z_L}.
\]
The exact expansion $R_L=m_L\sum_A\delta^{|A|}\chi_A$ proves that
$w_L$ has mean one and the required zero coefficients. Moreover
\[
 \|p_L\|_\infty\leq m_L\sum_{j=1}^{k_L}\binom Lj\delta^j
 \leq m_L k_LL^{k_L}\longrightarrow0.
\]
Indeed $\log m_L=-(s-1)d^2L^{1/4+\alpha/2}/2+O(L\delta^4)$,
and $\log k_L+k_L\log L=o(L^{1/4+\alpha/2})$.
It follows that $w_L\geq(1-c_L\|p_L\|_\infty)/Z_L\to1$ in this
lower-bound sense, so eventually $w_L\geq b$. Also
$\|w_L-v_L\|_\infty\to0$, $\|w_L\|_s\to B_*$, and
$\|S_{w_L}-S_{v_L}\|\leq\|w_L-v_L\|_\infty\to0$.
The uniform baseline, together with this sup-norm bound, supplies the
positivity that an unrestricted Fourier erasure would not preserve.

Finally let $B_L=\|w_L\|_s$, and set $\varepsilon_L=0$ when
$B_L\leq B_*$, otherwise
$\varepsilon_L=(B_L-B_*)/(B_L-1)$. On the latter event the denominator
is positive, $0<\varepsilon_L<1$, and $\varepsilon_L\to0$.
Then $r_L=(1-\varepsilon_L)w_L+\varepsilon_L$ retains the exact
marginals, mean, and eventual floor. Minkowski gives $\|r_L\|_s\leq B_*$.
Endpoint hypercontractivity and H\"older, with
$p=2s/(s-1)=1+1/\rho$, imply for every real signed $f$ that
\begin{equation}\label{eq:critical-signed-322}
 \|QM_fQ\|\leq\|f\|_s,
 \qquad \|Q\|_{2\to p}\leq1.
\end{equation}
Thus $\|S_{r_L}-S_{w_L}\|\leq\varepsilon_L(B_L+1)\to0$.
The actual leading eigenvalues still tend to one. Replace the finitely
many densities before all positivity and floor conditions hold by $r_L=1$.
Every stated constraint then holds in every dimension. The positive kernel
has a simple Perron root in each dimension, while its normalized gap tends
to zero. Below $B_*$ the earlier sign-compression bound
$2\lambda_2^s\leq\mathbb E r^s$ already supplies a uniform gap.
\end{proof}

\begin{theorem}[Rigidity from critical-power uniform integrability]
\label{thm:marginal-ui}
Suppose $k_L\to\infty$, $r_L\geq0$, $\mathbb E r_L=1$, and the
weighting laws have uniform marginals through order $k_L$. If
\[
 \lim_{M\to\infty}\sup_L
 \mathbb E[r_L^s\mathbf1_{r_L>M}]=0,
\]
then $\|S_{r_L}-P_\rho\|\to0$ and $\gap(T_{r_L})\to1-\rho$.
For $\rho<\sigma<1$, put $t=\sqrt{\rho/\sigma}$,
$m_L=\lfloor k_L/2\rfloor$, and
$\eta_M=\sup_L\|r_L\mathbf1_{r_L>M}\|_s$. More precisely,
\[
 \|S_{r_L}-P_\rho\|
 \leq2(M+1)t^{m_L+1}+2\eta_M.
\]
\end{theorem}
\begin{proof}
Let $\Pi$ project onto Walsh degrees at most $m_L$. With $h=r_L-1$,
the $(A,C)$ entry of $QM_hQ$ is
$\rho^{(|A|+|C|)/2}\widehat h(A\mathbin\triangle C)$.
Since $|A\mathbin\triangle C|\leq2m_L\leq k_L$ on this block,
including the zero constant coefficient, $\Pi QM_hQ\Pi=0$.

Split $h=h_b+h_t$, with $h_b=r_L\mathbf1_{r_L\leq M}-1$ and
$h_t=r_L\mathbf1_{r_L>M}$. Write $Q_\sigma=P_{\sqrt\sigma}$ and
$D_t=P_t$. For $B_b=QM_{h_b}Q$ we have
$B_b=D_t A_bD_t$, $A_b=Q_\sigma M_{h_b}Q_\sigma$,
$\|A_b\|\leq M+1$, and $\|D_t(I-\Pi)\|\leq t^{m_L+1}$.
The identity
$B_b-\Pi B_b\Pi=(I-\Pi)B_b+\Pi B_b(I-\Pi)$ bounds its norm by
$2(M+1)t^{m_L+1}$. For $B_t=QM_{h_t}Q$,
\eqref{eq:critical-signed-322} gives $\|B_t-\Pi B_t\Pi\|\leq2\eta_M$.
Since $\Pi(B_b+B_t)\Pi=0$, addition proves the estimate.
First send $L\to\infty$ at fixed $M$, then $M\to\infty$.
Self-adjoint perturbation and the actual spectrum $1,\rho,\ldots$ of
$P_\rho$ give the first two eigenvalue limits $1,\rho$ and the gap limit.
\end{proof}

For comparison, a strictly higher moment gives an explicit rate without
truncation. If $q>s$, $\|r_L\|_q\leq B$, and
$\rho<\sigma\leq(q-1)/(q+1)$, the same factorization gives
\[
 \|S_{r_L}-P_\rho\|\leq2(B+1)
 (\rho/\sigma)^{(\lfloor k_L/2\rfloor+1)/2}.
\]
Here $Q_\sigma M_{r_L-1}Q_\sigma$ has norm at most $B+1$ by
H\"older and hypercontractivity. The corresponding low block vanishes
exactly. This estimate and the preceding theorem are sufficient conditions,
not characterizations of all stable families.

\begin{theorem}[Rigidity for every positive linear marginal fraction]
\label{thm:linear-marginal-rigidity}
Suppose $r_L\geq0$, $\mathbb E r_L=1$, $\|r_L\|_s\leq B<\infty$,
and every marginal through order $k_L\leq L$ is uniform. If
$\liminf_L k_L/L>0$, then $\|S_{r_L}-P_\rho\|\to0$.
In particular the two leading eigenvalues of $T_{r_L}$ tend to $1,\rho$,
and $\gap(T_{r_L})\to1-\rho$. No density floor, uniform integrability,
or supercritical moment assumption is required.
\end{theorem}

\begin{proof}
Set $p=2s/(s-1)=1+1/\rho>2$. We use Theorem~1.3 and
Proposition~2.7 of the version-pinned primary source
\cite{kirshner-samorodnitsky}. In uniform probability normalization, its
continuous function $\psi(p,x)$ on $[0,1/2]$ bounds every real homogeneous
degree-$j$ Walsh polynomial, $j\leq L/2$, by
\[
 \|f_j\|_p\leq2^{L\psi(p,j/L)/p}\|f_j\|_2.
\]
For fixed $p$, $\psi(p,0)=0$, the right derivative at zero is
$(p/2)\log_2(p-1)$, and the function is strictly concave.
These are attributed external mathematical inputs. Strict concavity gives
\[
 \Delta_p(x)=\frac{x\log(p-1)}2-
 \frac{\log2}{p}\psi(p,x)>0\quad(0<x\leq1/2).
\]
For completeness, if $f(0)=0$ is strictly concave with right derivative
$d$, then $f(x)/x<f(x/2)/(x/2)\leq d$; this also handles the upper
endpoint. Natural logarithms are used except for the displayed $\log_2$.

For $x\in[0,1]$ let $y=\min(x,1-x)$ and
\[
 a_p(x)=\frac{x\log(p-1)}2-\frac{\log2}{p}\psi(p,y)
 =\Delta_p(y)+\frac{(x-y)\log(p-1)}2.
\]
This is continuous and strictly positive for every $x>0$, including
$a_p(1)=\log(p-1)/2$. For each fixed $0<\gamma\leq1$, compactness gives
$a_\gamma=\min_{x\in[\gamma,1]}a_p(x)>0$.
For $j\leq L/2$ the published inequality and $\rho=1/(p-1)$ imply
\[
 \|Qf_j\|_p\leq e^{-La_p(j/L)}\|f_j\|_2.
\]
For $j>L/2$, multiplication by the full parity character maps the entire
degree-$j$ level to degree $L-j$ while preserving all $L^q$ norms.
Applying the preceding homogeneous inequality there and retaining the
original factor $\rho^{j/2}$ gives the same bound. This does not assert
that parity multiplication commutes with $Q$.

Let $H_\gamma$ project onto degrees $j\geq\gamma L$. Triangle inequality,
then Cauchy--Schwarz over at most $L+1$ orthogonal levels, give
\begin{equation}\label{eq:high-level-endpoint-322}
 \|QH_\gamma\|_{2\to p}\leq\sqrt{L+1}\,e^{-La_\gamma}.
\end{equation}
Take fixed $0<c<\liminf k_L/L$ and put $\gamma=c/2$. For all sufficiently
large $L$, the complement $H=I-\Pi$ of degrees at most
$m_L=\lfloor k_L/2\rfloor$ has degrees
$j\geq m_L+1\geq k_L/2\geq\gamma L$.
For $\mathcal B=QM_{r_L-1}Q$, exact marginals give $\Pi\mathcal B\Pi=0$,
hence $\mathcal B=H\mathcal B+\Pi\mathcal BH$.
For real unit vectors $u,v$, H\"older with $1/s+2/p=1$ yields
\[
 |\langle u,H\mathcal Bv\rangle|
 \leq\|r_L-1\|_s\|QHu\|_p\|Qv\|_p
 \leq(B+1)\sqrt{L+1}\,e^{-La_\gamma}.
\]
The other term uses $Q\Pi u$ and $QHv$. Since $\Pi$ is an $L^2$
contraction, $\|Q\Pi\|_{2\to p}\leq1$; no $L^p$ contraction property
of $\Pi$ is needed. Thus
\[
 \|S_{r_L}-P_\rho\|\leq2(B+1)\sqrt{L+1}\,e^{-La_\gamma}.
\]
This tends to zero. The ordered-eigenvalue perturbation bound applies in
the same finite dimensional space for each $L$. It gives the actual
top-two limits, including when $r_L$ has zeros, and the common spectrum
of $S_{r_L},T_{r_L}$ supplies the asserted gap limit.
\end{proof}

\paragraph{Scope of the separation.}
Theorems~\ref{thm:growing-collapse} and~\ref{thm:linear-marginal-rigidity}
operate at the same critical cap $B_*$. Their conclusions do not locate a
sharp growth threshold between sub-square-root powers and linear order.
In particular qualitative strict concavity does not provide a uniform
quadratic deficit near zero, and the constant $a_\gamma$ above is fixed
only for fixed $\gamma>0$. No shrinking-cutoff rate is inferred.
Theorem~\ref{thm:marginal-ui} supplies a separate sufficient condition at
arbitrary diverging marginal order. All statements concern finite Boolean
transfer operators with global weights. They imply no fixed-range Gibbs,
local-dynamics, physical-gap, continuum, or Yang--Mills result.
The classical marginal/Fourier criterion and coefficient erasure are prior
methods; the homogeneous moment improvement is explicitly attributed.
Historical priority for their joint spectral application remains unestablished.


\section*{Model assessments and review context}
OpenAI/Codex assisted the mathematical synthesis and manuscript preparation.
The fabrication session is involved in the work. A distinct involved
Codex/OpenAI session audits the frozen mathematical source; exact served
backend versions and reasoning settings are not exposed for these source
audits, so no particular
model version or provider independence is asserted. Analytical source
review is not human peer review and supplies no inherited PDF rating.
The tail audits300--303 used one distinct auditor context with its earlier
readings retained, not four blind or provider-independent reviews. The
cutoff-one source306 received a separate fresh-assignment AI audit;
shared provider and filesystem do not establish technical isolation or
independence. The
historical v4 PDF review records a requested configuration separately from
what was actually exposed; its rating applies only to its original PDF.
Source309's first audit found a missing upper restriction on a floor
parameter in an additional comparison; that failed version remains
unchanged. Its corrected version was reread in the same involved auditor
context. Source310 received a separate fresh-assignment source audit;
shared provider and filesystem again preclude an independence claim.
The v5 exact-PDF review requested gpt-6-astra with high reasoning effort,
but the served backend and effort were not independently attested.
Its full report and six reservations remain bound to its original PDF,
not this revision. No numerical assessment of the present revision
has yet been incorporated or inherited.
Source313 received a separate fresh-assignment analytical audit. Its
reader checked the inherited explicit balanced family rather than relying
on previous acceptance labels. The shared provider and filesystem again
do not establish independence. Its minor request to bracket each
normalized product factor is incorporated here. The prior311 document
audit's notation and navigation observations are also addressed, with all
twenty-one prior proofs and twenty prior statement environments preserved.
None of these source checks is an exact-PDF assessment or human review.
The scalar-envelope and universal-tail source checks use a distinct
involved auditor context with prior source313 readings retained. They are
not blind or provider-independent reviews. The historical v6 PDF review
requested gpt-6-astra with high reasoning effort, with served backend and
effort not independently attested. Its full report remains bound to its
original PDF. The new mathematical claims require assessment on this
revision's exact final PDF and do not inherit its numerical rating.
The grouped growing-marginal construction and linear-fraction rigidity
argument received separate exact-source analytical audits in an involved
Codex/OpenAI context, with prior readings and a shared filesystem. The
Kirshner--Samorodnitsky inequality is a credited external input, not reproved
here. Full reports, fixed checks and source identities are retained with this
revision. The v7 exact-PDF assessment and all five reservations remain
historical evidence bound to PDF SHA-256
\texttt{ad708fafb3de9fc83cd2dc05443f34b2792da94a20d0bf1b1cd7b8ca41b9340c}.
Its page-22 reference to an unrelated equation is repaired here by a symbolic
reference to the mixed-power inequality. Neither the old numerical grade nor
an independent-review claim is transferred to this revised manuscript.
Any assessment attached to this manuscript must identify the exact final
PDF hash, its actual review context, complete report, scores and objections.
The manuscript does not incorporate a grade from another version or paper.
The author owns and administers AIRR. A working-paper deposit is not an
editorial acceptance or evidence of human peer review.

\smallskip\noindent
Reproducible source commit:\par\noindent
{\small\ttfamily\SourceCommit}

\begin{thebibliography}{99}
\bibitem{BertsimasPopescu}
D. Bertsimas and I. Popescu,
Optimal Inequalities in Probability Theory: A Convex Optimization Approach,
\emph{SIAM Journal on Optimization} \textbf{15}(3) (2005), 780--804.
\href{https://doi.org/10.1137/S1052623401399903}{Publisher version}.
\bibitem{Aida}
S. Aida, An estimate of the gap of spectrum of Schr\"odinger operators
which generate hyperbounded semigroup, author preprint, 9 February 2001.
\href{https://www.ms.u-tokyo.ac.jp/~aida/paper/preprint/handai3.pdf}{Author preprint}.
\bibitem{BiskupKoenig}
M. Biskup and W. K\"onig, Eigenvalue order statistics for random
Schr\"odinger operators with doubly-exponential tails,
author manuscript, 21 April 2015, Lemma~4.5 and Eq.~(4.29).
\href{https://www.wias-berlin.de/people/koenig/www/PPP4PAM-CMP-final.pdf}{Author manuscript}.
\bibitem{AGM}
N. Alon, O. Goldreich and Y. Mansour,
Almost $k$-wise independence versus $k$-wise independence,
\emph{Information Processing Letters} \textbf{88} (2003), 107--110.
\href{https://www.wisdom.weizmann.ac.il/~oded/PSX/agm.pdf}{Author manuscript}.
\bibitem{BrehmResch}
M. Brehm and N. Resch,
Linear Time Encodable Binary Code Achieving GV Bound with Linear Time
Encodable Dual Achieving GV Bound,
\emph{ITCS 2026}, LIPIcs \textbf{362}, article 28.
\href{https://doi.org/10.4230/LIPIcs.ITCS.2026.28}{Publisher version}.
\bibitem{ErikssonCongruence}
L. Eriksson,
Congruence Rigidity and the Fusion Bound: What a Positive Weight Can and
Cannot do to the Spectrum of a Transfer Kernel,
AIRR, ARR-2026-7S1XA3PZEB9QM8FQ, version 1.
\href{https://airr.science/papers/ARR-2026-7S1XA3PZEB9QM8FQ/versions/v1/}{Public record}.
\bibitem{ErikssonParity}
L. Eriksson,
Strictly Positive Parity Laws and the Limits of Marginal Certificates:
Exact Heat-Bath Spectrum and Optimal Variance Tensorization,
AIRR, ARR-2026-5HRC3PS9HZ9WDAA5, version 2 (2026).
\href{https://airr.science/papers/ARR-2026-5HRC3PS9HZ9WDAA5/versions/v2/}{Public working paper}.
\bibitem{MORSS}
E. Mossel, R. O'Donnell, O. Regev, J. E. Steif and B. Sudakov,
Non-interactive correlation distillation, inhomogeneous Markov chains,
and the reverse Bonami-Beckner inequality,
\emph{Israel Journal of Mathematics} \textbf{154} (2006), 299--336.
\href{https://www.cs.cmu.edu/~odonnell/papers/nicd.pdf}{Author manuscript}.
\bibitem{Cattiaux}
P. Cattiaux, Hypercontractivity for perturbed diffusion semigroups,
\emph{Annales de la Facult\'e des Sciences de Toulouse}
\textbf{14}(4) (2005), 609--628.
\href{https://www.numdam.org/item/AFST_2005_6_14_4_609_0/}{Publisher archive}.
\bibitem{MOS}
E. Mossel, K. Oleszkiewicz and A. Sen,
On Reverse Hypercontractivity, arXiv:1108.1210v5 (2012).
\href{https://arxiv.org/abs/1108.1210}{Primary manuscript}.
\bibitem{Gross}
L. Gross, Invariance of intrinsic hypercontractivity under perturbation of
Schr\"odinger operators, arXiv:2412.20282v2 (2025).
\href{https://arxiv.org/abs/2412.20282}{Primary manuscript}.
\bibitem{Gluck}
J. Gl\"uck, Spectral gaps for hyperbounded operators,
arXiv:1802.09422v6 (2019).
\href{https://arxiv.org/abs/1802.09422}{Primary manuscript}.
\bibitem{Miclo}
L. Miclo, On hyperboundedness and spectrum of Markov operators.
\href{https://perso.math.univ-toulouse.fr/miclo/files/2012/04/hyperboundedness.pdf}{Author manuscript}.
\bibitem{Wang}
F.-Y. Wang, Criteria of Spectral Gap for Markov Operators,
arXiv:1305.4460v6 (2013).
\href{https://arxiv.org/abs/1305.4460v6}{Primary manuscript}.
\bibitem{kirshner-samorodnitsky}
N. Kirshner and A. Samorodnitsky,
A moment ratio bound for polynomials and some extremal properties of
Krawchouk polynomials and Hamming spheres,
arXiv:1909.11929v1 (2019), Theorem~1.3 and Proposition~2.7.
\href{https://arxiv.org/abs/1909.11929v1}{Version-pinned primary source}.
\bibitem{alon-kwise}
N. Alon, A. Andoni, T. Kaufman, K. Matulef, R. Rubinfeld and N. Xie,
Testing k-wise and Almost k-wise Independence,
\emph{Proceedings of STOC} (2007).
\href{https://www.mit.edu/~andoni/papers/kwise.pdf}{Author manuscript},
Fact~2.3 and Section~3.
\bibitem{rubinfeld-xie}
R. Rubinfeld and N. Xie,
Testing Non-uniform k-wise Independent Distributions over Product Spaces,
author manuscript, Appendix~A.2.
\href{https://people.csail.mit.edu/ningxie/papers/RX09.pdf}{Primary source}.
\end{thebibliography}
\end{document}